%======================================================
\documentclass[12pt]{article}
%======================================================
\usepackage[utf8]{inputenc}
\usepackage[T1]{fontenc}
\usepackage{enumitem}
\usepackage[spanish,mexico]{babel}
\usepackage{lmodern}
\usepackage{amsmath,amssymb,amsthm,mathtools}
\usepackage{geometry,enumitem,booktabs,array,xcolor,graphicx,hyperref,fancyhdr,microtype,longtable,multicol}
%======================================================
\geometry{letterpaper,top=2.4cm,bottom=2.35cm,left=2.45cm,right=2.45cm}
\setlength{\parindent}{0pt}
\setlength{\parskip}{0.52em}
\setlength{\headheight}{15pt}
\setlength{\emergencystretch}{2em}
%======================================================
\definecolor{azul}{RGB}{22,72,120}
\hypersetup{colorlinks=true,linkcolor=azul,urlcolor=azul,pdftitle={Algebra Lineal - Unidad 3: Matrices y determinantes},pdfauthor={Carlos Ernesto Martinez}}
%======================================================
\pagestyle{fancy}
\fancyhf{}
\lhead{\small Tecnológico Nacional de México}
\rhead{\small Álgebra Lineal}
\lfoot{\footnotesize Carlos Ernesto Martínez}
\cfoot{\footnotesize\thepage}
\rfoot{\footnotesize mrdgzcarlose@gmail.com}
%======================================================
\renewcommand{\headrulewidth}{0.4pt}
\renewcommand{\footrulewidth}{0.3pt}
\newtheorem{definicion}{Definición}[section]
\newtheorem{teorema}[definicion]{Teorema}
\newtheorem{proposicion}[definicion]{Proposición}
\newtheorem{corolario}[definicion]{Corolario}
\theoremstyle{remark}
\newtheorem{observacion}[definicion]{Observación}
\newtheorem{ejemplo}[definicion]{Ejemplo}
\newtheorem{Ejer}[definicion]{Ejercicio}
\newcommand{\R}{\mathbb{R}}
\newcommand{\C}{\mathbb{C}}
\newcommand{\Mat}[2]{\mathcal{M}_{#1\times #2}}
\newcommand{\rang}{\operatorname{rang}}
\newcommand{\nul}{\operatorname{nul}}
\newcommand{\tr}{\operatorname{tr}}
\newcommand{\adj}{\operatorname{adj}}
\newcommand{\Ker}{\operatorname{Ker}}
\newcommand{\RREF}{\operatorname{RREF}}
\newcommand{\Span}{\operatorname{span}}
\newcommand{\vect}[1]{\boldsymbol{#1}}
\newcommand{\idea}[1]{\begin{center}\fcolorbox{azul}{white}
%======================================================
{\parbox{0.91\linewidth}{#1}}\end{center}}
\newcounter{contadorEjercicio}
\newenvironment{listaejercicios}{\begin{enumerate}[label=\textbf{E\arabic*.},leftmargin=2.4em,itemsep=0.52em]\setcounter{enumi}{\value{contadorEjercicio}}}{\setcounter{contadorEjercicio}{\value{enumi}}\end{enumerate}}
%======================================================
\title{\textbf{Notas para el curso de Álgebra Lineal}\\[0.35cm]
\Large Unidad 3: Matrices y determinantes}
\author{Carlos Ernesto Martínez\\Tecnológico Nacional de México}
\date{Agosto--Diciembre 2026}
%======================================================
\begin{document}
\maketitle
%\begin{center}\textcolor{azul}{\rule{0.72\textwidth}{0.6pt}}\end{center}
%======================================================
\tableofcontents
\newpage
%=======================================================
\section*{Propósito y organización de la unidad}
%=======================================================
\textbf{Contenidos del temario.} Esta unidad cubre los 27 subtemas de la Unidad 3: definición y clasificación de matrices, operaciones, transposición, inversa, determinantes, menores, cofactores, adjunta, rango, núcleo y aplicaciones matriciales.

%=========================================================
\section{Matrices: definición, orden y clasificación}
%==========================================================

Recordemos la expresión matricial del sistema de ecuaciones lineales:
\[
\left\{\begin{aligned}2x-y+3z&=4,\\x+2y-z&=1\end{aligned}\right.
\quad\Longleftrightarrow\quad
\underbrace{\begin{pmatrix}2&-1&3\\1&2&-1\end{pmatrix}}_{A\in\Mat{2}{3}}
\underbrace{\begin{pmatrix}x\\y\\z\end{pmatrix}}_{\vect x\in\R^3}
=\underbrace{\begin{pmatrix}4\\1\end{pmatrix}}_{\vect b\in\R^2}.
\]
La expresión $A\vect x=\vect b$ exige que el número de columnas de $A$ coincida con el número de componentes de $\vect x$. El vector resultante tiene tantas componentes como filas tiene $A$.

\[
A=\begin{pmatrix}|&|&|\\\vect a_1&\vect a_2&\vect a_3\\|&|&|\end{pmatrix},\qquad
A\vect x=x_1\vect a_1+x_2\vect a_2+x_3\vect a_3.
\]


\begin{proposicion}
Para demostrar que dos matrices del mismo orden son iguales es suficiente demostrar que coinciden sus entradas de posición $(i,j)$ para cada par válido $(i,j)$.
\end{proposicion}
%======================================================
\subsection*{Definición y notación}
%======================================================
Una matriz permite organizar datos o coeficientes mediante filas y columnas. El objeto resultante no es sólo una tabla: su estructura hace posible definir operaciones algebraicas con propiedades precisas.

\begin{definicion}[Matriz]
Una \emph{matriz real de orden $m\times n$} es un arreglo rectangular de $mn$ números reales, distribuidos en $m$ filas y $n$ columnas:
\[
A=\begin{pmatrix}
a_{11}&a_{12}&\cdots&a_{1n}\\
a_{21}&a_{22}&\cdots&a_{2n}\\
\vdots&\vdots&\ddots&\vdots\\
a_{m1}&a_{m2}&\cdots&a_{mn}
\end{pmatrix}=\bigl(a_{ij}\bigr)_{m\times n}.
\]
El número $a_{ij}$ es la entrada situada en la fila $i$ y la columna $j$. El conjunto de todas las matrices reales $m\times n$ se denota por $\Mat{m}{n}(\R)$.
\end{definicion}

\begin{ejemplo}
Sea
\[
A=\begin{pmatrix}2&-1&3\\0&4&5\end{pmatrix}.
\]
La matriz $A$ tiene dos filas y tres columnas, así que su orden es $2\times3$. \end{ejemplo}

\begin{definicion}
Dos matrices $A=(a_{ij})$ y $B=(b_{ij})$ son iguales si y sólo si tienen el mismo orden y sus entradas correspondientes coinciden:
\[
A=B\quad\Longleftrightarrow\quad a_{ij}=b_{ij}\quad\text{para todo }i,j.
\]
\end{definicion}

\begin{ejemplo}
Encontrar $x,y,z$ si
\[
\begin{pmatrix}x+1&2y\\z-3&5\end{pmatrix}=
\begin{pmatrix}4&-6\\1&5\end{pmatrix}.
\]
Igualamos las entradas correspondientes:
\[
x+1=4,\qquad 2y=-6,\qquad z-3=1.
\]
Por tanto, $\boxed{x=3,\ y=-3,\ z=4}$.
\end{ejemplo}
%=======================================================
\subsection*{Tipos de matrices}
%=======================================================
\begin{definicion}[Matrices especiales]
Consideremos una matriz $A=(a_{ij})$.
\begin{itemize}
\item \textbf{Cuadrada:} tiene el mismo número de filas y columnas ($n\times n$).
\item \textbf{Nula:} todas sus entradas son cero; se denota por $0$ cuando el orden está claro.
\item \textbf{Diagonal:} es cuadrada y $a_{ij}=0$ cuando $i\ne j$.
\item \textbf{Identidad:} es diagonal, con $a_{ii}=1$ para todo $i$; se denota por $I_n$.
\item \textbf{Triangular superior:} $a_{ij}=0$ si $i>j$.
\item \textbf{Triangular inferior:} $a_{ij}=0$ si $i<j$.
\item \textbf{Simétrica:} $A=A^T$.
\end{itemize}
\end{definicion}

\begin{observacion}
Una matriz diagonal es siempre simétrica.
\end{observacion}

%=======================================================

Llamamos \emph{escalar} a una matriz de la forma $\lambda I_n$. Una matriz \emph{antisimétrica} cumple $A^T=-A$, lo que, trabajando sobre $\R$, obliga a que sus entradas diagonales sean cero.
\begin{proposicion}
Toda matriz cuadrada real $A$ admite la descomposición única
\[
\boxed{A=S+K,\qquad S=\tfrac12(A+A^T),\quad K=\tfrac12(A-A^T),}
\]
donde $S$ es simétrica y $K$ es antisimétrica.
\end{proposicion}
\begin{proof}
Por las propiedades de la transposición,
\[
S^T=\tfrac12(A^T+A)=S,\qquad K^T=\tfrac12(A^T-A)=-K.
\]
Sumando las dos expresiones obtenemos $S+K=A$. Si $A=S_1+K_1$ con $S_1^T=S_1$ y $K_1^T=-K_1$, entonces $A^T=S_1-K_1$; de allí $S_1=(A+A^T)/2$ y $K_1=(A-A^T)/2$. Esto también demuestra unicidad.
\end{proof}
\begin{ejemplo}
Para $A=\begin{psmallmatrix}1&4\\2&3\end{psmallmatrix}$,
\[
S=\frac12\begin{pmatrix}2&6\\6&6\end{pmatrix}=\begin{pmatrix}1&3\\3&3\end{pmatrix},\quad
K=\frac12\begin{pmatrix}0&2\\-2&0\end{pmatrix}=\begin{pmatrix}0&1\\-1&0\end{pmatrix}.
\]
Compruébese que $S+K=A$, $S^T=S$ y $K^T=-K$.
\end{ejemplo}

%=======================================================
\section{Operaciones con matrices}
%=======================================================

%=======================================================
\subsection*{Suma y multiplicación por un escalar}
%=======================================================
\begin{definicion}
Para matrices $A=(a_{ij})$ y $B=(b_{ij})$ del \emph{mismo orden}, definimos
\[
A+B=(a_{ij}+b_{ij}).
\]
La suma no está definida, bajo esta operación, si los órdenes de $A$ y $B$ son diferentes.
\end{definicion}

\begin{definicion}
Sea $\lambda\in\R$ y $A=(a_{ij})$. Definimos
\[
\lambda A=(\lambda a_{ij}).
\]
El orden de $\lambda A$ es el mismo de $A$.
\end{definicion}

\begin{ejemplo}
Sean $A=\begin{psmallmatrix}1&-2\\3&0\end{psmallmatrix}$ y $B=\begin{psmallmatrix}4&1\\-1&5\end{psmallmatrix}$. Entonces
\[
A+B=\begin{pmatrix}5&-1\\2&5\end{pmatrix},\qquad
2A-B=\begin{pmatrix}-2&-5\\7&-5\end{pmatrix}.
\]
\end{ejemplo}

\begin{proposicion}
Para $A,B,C\in\Mat{m}{n}(\R)$ y $\lambda,\mu\in\R$ se cumple
\begin{eqnarray*}
\begin{array}{ll}
A+B=B+A,& (A+B)+C=A+(B+C),\\
A+0=A,& A+(-A)=0,\\
\lambda(A+B)=\lambda A+\lambda B,& (\lambda+\mu)A=\lambda A+\mu A,\\
(\lambda\mu)A=\lambda(\mu A),&1A=A.
\end{array}
\end{eqnarray*}

\end{proposicion}

Las matrices $m\times n$ se suman y multiplican por números reales exactamente como vectores de $\R^{mn}$. Podemos verificar \emph{todos} los axiomas de espacio vectorial mediante sus entradas. Por ejemplo, la asociatividad se deduce de
\[
((A+B)+C)_{ij}=(a_{ij}+b_{ij})+c_{ij}=a_{ij}+(b_{ij}+c_{ij})=(A+(B+C))_{ij}.
\]
La existencia del inverso aditivo se verifica con $(-A)_{ij}=-a_{ij}$ y la distributividad con
\[
\bigl(\lambda(A+B)\bigr)_{ij}=\lambda(a_{ij}+b_{ij})=\lambda a_{ij}+\lambda b_{ij}.
\]

\begin{ejemplo}
Sean
\[
A=\begin{pmatrix}1&0&-2\\3&1&4\\0&-1&2\end{pmatrix},\qquad
B=\begin{pmatrix}0&2&1\\-1&3&0\\5&-2&1\end{pmatrix}.
\]
Para obtener $2A-3B$ formamos primero
\[
2A=\begin{pmatrix}2&0&-4\\6&2&8\\0&-2&4\end{pmatrix},\qquad
3B=\begin{pmatrix}0&6&3\\-3&9&0\\15&-6&3\end{pmatrix}.
\]
La resta, realizada \emph{posición por posición}, produce
\[
\boxed{2A-3B=\begin{pmatrix}2&-6&-7\\9&-7&8\\-15&4&1\end{pmatrix}.}
\]
\end{ejemplo}
%=======================================================
\subsection*{Producto de matrices}
%=======================================================
El producto matricial no se calcula multiplicando simplemente entradas que ocupen la misma posición. La regla combina las filas de la primera matriz con las columnas de la segunda.

\begin{definicion}[Producto matricial]
Sean $A=(a_{ij})\in\Mat{m}{n}(\R)$ y $B=(b_{jk})\in\Mat{n}{p}(\R)$. Se define $C=AB\in\Mat{m}{p}(\R)$ mediante
\[
\boxed{\quad c_{ik}=\sum_{j=1}^{n}a_{ij}b_{jk}\quad}
\]
para $1\le i\le m$, $1\le k\le p$.
\end{definicion}

La condición de existencia del producto es que el número de columnas de $A$ coincida con el número de filas de $B$:
\[
\underbrace{(m\times n)}_A\ \underbrace{(n\times p)}_B
\longrightarrow\underbrace{(m\times p)}_{AB}.
\]

\begin{ejemplo}[Producto de una matriz $2\times3$ por una $3\times2$]
Calculemos $AB$ para
\[
A=\begin{pmatrix}1&2&-1\\0&3&4\end{pmatrix},\qquad
B=\begin{pmatrix}2&1\\-1&0\\3&2\end{pmatrix}.
\]
El producto existe y tiene orden $2\times2$. Calculamos cada entrada:
\begin{align*}
c_{11}&=1(2)+2(-1)+(-1)(3)=-3,\\
c_{12}&=1(1)+2(0)+(-1)(2)=-1,\\
c_{21}&=0(2)+3(-1)+4(3)=9,\\
c_{22}&=0(1)+3(0)+4(2)=8.
\end{align*}
Por tanto,
\[\boxed{AB=\begin{pmatrix}-3&-1\\9&8\end{pmatrix}.}\]
\end{ejemplo}

\begin{observacion}
Aunque el producto $AB$ exista, $BA$ podría no existir. Cuando ambos existen, generalmente tienen órdenes distintos o no son iguales.
\end{observacion}

\begin{ejemplo}[El producto no es conmutativo]
Sean
\[
A=\begin{pmatrix}1&1\\0&1\end{pmatrix},\qquad
B=\begin{pmatrix}1&0\\1&1\end{pmatrix}.
\]
Entonces
\[
AB=\begin{pmatrix}2&1\\1&1\end{pmatrix},\qquad
BA=\begin{pmatrix}1&1\\1&2\end{pmatrix}.
\]
Ambos productos existen, pero $AB\ne BA$. En consecuencia, \emph{no} podemos usar indiscriminadamente las reglas conmutativas de los números reales.
\end{ejemplo}

\begin{proposicion}[Propiedades del producto]
Si las dimensiones permiten realizar las operaciones, entonces
\begin{align*}
A(BC)&=(AB)C,\\
A(B+C)&=AB+AC,\\
(A+B)C&=AC+BC,\\
(\lambda A)B&=\lambda(AB)=A(\lambda B),\\
I_mA&=A=AI_n,\qquad A\in\Mat{m}{n}(\R).
\end{align*}
\end{proposicion}
\begin{proof}[Demostración de la asociatividad]
Sean $A$ de orden $m\times n$, $B$ de orden $n\times p$ y $C$ de orden $p\times q$. La entrada $(i,k)$ de $(AB)C$ es
\[
\sum_{\ell=1}^{p}\left(\sum_{j=1}^{n}a_{ij}b_{j\ell}\right)c_{\ell k}
=\sum_{j=1}^{n}\sum_{\ell=1}^{p}a_{ij}b_{j\ell}c_{\ell k}.
\]
La entrada $(i,k)$ de $A(BC)$ es
\[
\sum_{j=1}^{n}a_{ij}\left(\sum_{\ell=1}^{p}b_{j\ell}c_{\ell k}\right)
=\sum_{j=1}^{n}\sum_{\ell=1}^{p}a_{ij}b_{j\ell}c_{\ell k}.
\]
Son iguales porque las sumas son finitas; por tanto, $(AB)C=A(BC)$.
\end{proof}
%=======================================================
\subsection*{Interpretación del producto como combinación lineal de columnas}
%=======================================================
Si $A=[\vect a_1\ \cdots\ \vect a_n]$ y $\vect x=(x_1,\ldots,x_n)^T$, entonces la entrada $i$-ésima de $A\vect x$ es $\sum_{j=1}^na_{ij}x_j$. Por consiguiente,
\[
\boxed{A\vect x=x_1\vect a_1+\cdots+x_n\vect a_n.}
\]
Esta fórmula explica por qué cada columna $j$ de $AB$ se obtiene multiplicando $A$ por la columna $j$ de $B$.
\begin{ejemplo}[Producto $3\times3$ por columnas]
\[
A=\begin{pmatrix}1&2&0\\0&1&-1\\2&0&1\end{pmatrix},\quad
B=\begin{pmatrix}1&0&2\\0&1&1\\1&-1&0\end{pmatrix}.
\]
La primera columna del producto es $A(1,0,1)^T=(1,-1,3)^T$; la segunda es $A(0,1,-1)^T=(2,2,-1)^T$; la tercera es $A(2,1,0)^T=(4,1,4)^T$. Por tanto,
\[
\boxed{AB=\begin{pmatrix}1&2&4\\-1&2&1\\3&-1&4\end{pmatrix}.}
\]
Recomendación de cálculo: indicar primero las dimensiones, después obtener las entradas por productos fila--columna y, finalmente, verificar al menos una columna mediante combinaciones lineales.
\end{ejemplo}

%=======================================================
\subsection*{Propiedades que NO se heredan de los números reales}
%=======================================================
\begin{proposicion}
Existen matrices cuadradas no nulas $A$ y $B$ tales que $AB=0$.
\end{proposicion}
\begin{proof}
Elijamos $A=\begin{psmallmatrix}1&0\\0&0\end{psmallmatrix}$ y $B=\begin{psmallmatrix}0&0\\0&1\end{psmallmatrix}$. Por multiplicación directa, $AB=\begin{psmallmatrix}0&0\\0&0\end{psmallmatrix}$, aunque ninguna de las dos matrices es nula.
\end{proof}
La cancelación $AB=AC\implies B=C$ es válida cuando $A$ es invertible, pero no en general. En efecto, si existe $A^{-1}$, multiplicamos a la izquierda por $A^{-1}$ para obtener $B=C$.
\begin{proposicion}
Si $A\in\Mat{m}{n}(\R)$ y $B,C\in\Mat{n}{p}(\R)$, entonces $A(B+C)=AB+AC$.
\end{proposicion}
\begin{proof}
La entrada $(i,k)$ del lado izquierdo es $\sum_{j=1}^{n}a_{ij}(b_{jk}+c_{jk})$. La distributividad real permite separarla como $\sum_{j=1}^{n}a_{ij}b_{jk}+\sum_{j=1}^{n}a_{ij}c_{jk}$, que es exactamente la entrada $(i,k)$ de $AB+AC$.
\end{proof}

%=======================================================
\subsection*{Transpuesta de una matriz}
%=======================================================
\begin{definicion}
Para $A=(a_{ij})\in\Mat{m}{n}(\R)$, su \emph{transpuesta} es la matriz $A^T\in\Mat{n}{m}(\R)$ definida por
\[
(A^T)_{ji}=a_{ij}.
\]
En otras palabras, las filas se convierten en columnas.
\end{definicion}
\begin{ejemplo}
Si $A=\begin{psmallmatrix}1&-2&3\\4&0&5\end{psmallmatrix}$, entonces
\[
A^T=\begin{pmatrix}1&4\\-2&0\\3&5\end{pmatrix}.
\]
\end{ejemplo}

\begin{proposicion}
Para matrices de tamaños compatibles y $\lambda\in\R$:
\[
(A^T)^T=A,\qquad(A+B)^T=A^T+B^T,\qquad
(\lambda A)^T=\lambda A^T,\qquad (AB)^T=B^TA^T.
\]
\end{proposicion}
\begin{proof}[Demostración de $(AB)^T=B^TA^T$]
Sean $A\in\Mat{m}{n}(\R)$ y $B\in\Mat{n}{p}(\R)$. Para cada $i,k$:
\[
\bigl((AB)^T\bigr)_{ki}=(AB)_{ik}
=\sum_{j=1}^{n}a_{ij}b_{jk}
=\sum_{j=1}^{n}(B^T)_{kj}(A^T)_{ji}
=(B^TA^T)_{ki}.
\]
Por igualdad de entradas, la identidad queda demostrada.
\end{proof}

\begin{observacion}
Para matrices con entradas complejas, además de la transpuesta se utiliza la \emph{transpuesta conjugada} $A^*=\overline{A}^{\,T}$. La distinción será útil posteriormente al hablar de producto interno. 
\end{observacion}

%===============================================================
\section{Matrices inversas y el método de Gauss--Jordan}
%===============================================================
%=======================================================
\subsection*{Concepto de inversa}
%=======================================================
\begin{definicion}
Una matriz cuadrada $A\in\Mat{n}{n}(\R)$ es \emph{invertible} si existe una matriz $B$ del mismo orden tal que
\[
AB=BA=I_n.
\]
La matriz $B$ se llama \emph{inversa} de $A$ y se escribe $A^{-1}$. Si no existe, $A$ es \emph{singular} o \emph{no invertible}.
\end{definicion}

\begin{proposicion}
Si $A$ es invertible, su inversa es única.
\end{proposicion}
\begin{proof}
Supongamos que $B$ y $C$ son inversas de $A$. Usando asociatividad,
\[
B=BI_n=B(AC)=(BA)C=I_nC=C.
\]
Por tanto, no pueden existir dos inversas distintas.
\end{proof}

\begin{proposicion}
Si $A$ y $B$ son matrices cuadradas invertibles del mismo orden, entonces
\[
\boxed{(AB)^{-1}=B^{-1}A^{-1}.}
\]
\end{proposicion}
\begin{proof}
Basta comprobar
\[
(AB)(B^{-1}A^{-1})=A(BB^{-1})A^{-1}=AA^{-1}=I_n,
\]
y, análogamente, $(B^{-1}A^{-1})(AB)=I_n$. La unicidad de la inversa concluye la demostración.
\end{proof}
%=======================================================
\subsection*{Obtención de la inversa mediante Gauss--Jordan}
%=======================================================
La idea consiste en aplicar operaciones elementales simultáneamente a $A$ y a la identidad, mediante la matriz aumentada
\[
\left[A\mid I_n\right].
\]
Si las operaciones convierten el bloque izquierdo en $I_n$, el bloque derecho resulta ser $A^{-1}$:
\[
\boxed{[A\mid I_n]\ \sim\ [I_n\mid A^{-1}].}
\]
Cuando el bloque izquierdo no puede reducirse a la identidad, $A$ no es invertible.

\begin{ejemplo}
Sea
\[
A=\begin{pmatrix}2&1\\1&1\end{pmatrix}.
\]
Formamos la matriz aumentada y efectuamos las operaciones indicadas:
\begin{align*}
\left(\begin{array}{cc|cc}2&1&1&0\\1&1&0&1\end{array}\right)
&\xrightarrow{R_1\leftrightarrow R_2}
\left(\begin{array}{cc|cc}1&1&0&1\\2&1&1&0\end{array}\right)[1ex]
\xrightarrow{R_2\leftarrow R_2-2R_1}
\left(\begin{array}{cc|cc}1&1&0&1\\0&-1&1&-2\end{array}\right)\\[1ex]
&\xrightarrow{R_2\leftarrow -R_2}
\left(\begin{array}{cc|cc}1&1&0&1\\0&1&-1&2\end{array}\right)[1ex]
\xrightarrow{R_1\leftarrow R_1-R_2}
\left(\begin{array}{cc|cc}1&0&1&-1\\0&1&-1&2\end{array}\right).
\end{align*}
La matriz inversa es la que aparece en el bloque derecho:
\[
\boxed{A^{-1}=\begin{pmatrix}1&-1\\-1&2\end{pmatrix}.}
\]
Verificamos multiplicando:
\[
AA^{-1}=\begin{pmatrix}2&1\\1&1\end{pmatrix}
\begin{pmatrix}1&-1\\-1&2\end{pmatrix}
=\begin{pmatrix}1&0\\0&1\end{pmatrix}.
\]
\end{ejemplo}

\begin{observacion}
Resolver $A\vect{x}=\vect{b}$ mediante $\vect{x}=A^{-1}\vect{b}$ exige que $A$ sea cuadrada e invertible. En problemas grandes suele ser preferible resolver directamente el sistema sin construir explícitamente $A^{-1}$.
\end{observacion}

%===============================================================
\subsection*{Operaciones elementales y matrices elementales}
%=======================================================
La receta $[A\mid I_n]\sim[I_n\mid A^{-1}]$ tiene una justificación algebraica. Una operación elemental sobre filas equivale a multiplicar por la izquierda por una matriz \emph{elemental} $E$, obtenida al efectuar la misma operación sobre $I_n$.

\begin{definicion}
Son tres: $(i)$ intercambiar dos filas, $(ii)$ multiplicar una fila por un escalar no nulo y $(iii)$ sumar a una fila un múltiplo de otra. Cada operación es reversible; por ello su matriz elemental $E$ es invertible.
\end{definicion}

\begin{ejemplo}
Si $E=\begin{psmallmatrix}1&0&0\\-2&1&0\\0&0&1\end{psmallmatrix}$, entonces $EA$ coincide con $A$ después de aplicar $R_2\leftarrow R_2-2R_1$. Para invertir esta operación basta $E^{-1}=\begin{psmallmatrix}1&0&0\\2&1&0\\0&0&1\end{psmallmatrix}$.
\end{ejemplo}


\begin{teorema}
Sea $A$ cuadrada de orden $n$. Si una sucesión de operaciones elementales produce $E_s\cdots E_2E_1A=I_n$, entonces $A$ es invertible y
\[
\boxed{A^{-1}=E_s\cdots E_2E_1.}
\]
Por esta razón aplicar dichas operaciones a $[A\mid I_n]$ produce $[I_n\mid A^{-1}]$.
\end{teorema}
\begin{proof}
Sea $P=E_s\cdots E_1$. Todas las matrices elementales son invertibles, de modo que $P$ lo es. De $PA=I_n$ se sigue $A=P^{-1}$; por tanto, $AP=I_n$ y $P=A^{-1}$. Sobre el bloque derecho tenemos $PI_n=P=A^{-1}$.
\end{proof}

\begin{ejemplo}
Para $A=\begin{psmallmatrix}1&1&0\\0&1&1\\0&0&1\end{psmallmatrix}$, construimos
\[
\left(\begin{array}{ccc|ccc}1&1&0&1&0&0\\0&1&1&0&1&0\\0&0&1&0&0&1\end{array}\right).
\]
El pivote de la tercera columna ya es uno. Eliminamos primero la entrada encima de ese pivote:
\[
\xrightarrow{R_2\leftarrow R_2-R_3}
\left(\begin{array}{ccc|ccc}1&1&0&1&0&0\\0&1&0&0&1&-1\\0&0&1&0&0&1\end{array}\right).
\]
Eliminamos finalmente la entrada de la primera fila en la segunda columna:
\[
\xrightarrow{R_1\leftarrow R_1-R_2}
\left(\begin{array}{ccc|ccc}1&0&0&1&-1&1\\0&1&0&0&1&-1\\0&0&1&0&0&1\end{array}\right).
\]
La inversa es $A^{-1}=\begin{psmallmatrix}1&-1&1\\0&1&-1\\0&0&1\end{psmallmatrix}$. La verificación puede hacerse multiplicando $AA^{-1}$ y $A^{-1}A$.
\end{ejemplo}
%=======================================================
\subsection*{Condiciones para una inversa: una demostración operativa}
%=======================================================
\begin{teorema}
Para $A\in\Mat{n}{n}(\R)$ son equivalentes: (a) $A$ tiene inversa; (b) la forma reducida por filas es $I_n$; (c) $A\vect x=\vect0$ posee únicamente la solución trivial.
\end{teorema}
\begin{proof}
(a)$\Rightarrow$(c): si $A\vect x=\vect0$, multiplicar por $A^{-1}$ da $\vect x=\vect0$. (c)$\Rightarrow$(b): al llevar $A$ a forma reducida, cualquier columna sin pivote produciría una variable libre y una solución homogénea no trivial. Así que hay $n$ pivotes y la forma reducida es $I_n$. (b)$\Rightarrow$(a): el teorema de las matrices elementales proporciona una inversa.
\end{proof}

%=======================================================
\section{Determinantes}
%===============================================================
%=======================================================
\subsection*{Determinantes de orden dos y tres}
%=======================================================
El determinante asocia a cada matriz cuadrada un número. Este número permite distinguir, entre otras cosas, matrices invertibles de matrices singulares.

\begin{definicion}
Para $A=\begin{psmallmatrix}a&b\\c&d\end{psmallmatrix}$ definimos
\[
\det(A)=\begin{vmatrix}a&b\\c&d\end{vmatrix}=ad-bc.
\]
\end{definicion}
\begin{ejemplo}
Para $A=\begin{psmallmatrix}2&1\\1&1\end{psmallmatrix}$ tenemos
\[
\det(A)=2(1)-1(1)=1.
\]
El determinante no nulo anticipa que $A$ es invertible, como ya comprobamos por Gauss--Jordan.
\end{ejemplo}

\begin{definicion}
Si
\[
A=\begin{pmatrix}a&b&c\\d&e&f\\g&h&i\end{pmatrix},
\]
su determinante es
\[
\boxed{\det(A)=aei+bfg+cdh-ceg-bdi-afh.}
\]
Esta expresión puede organizarse mediante la llamada \emph{regla de Sarrus}, válida únicamente para matrices $3\times3$.
\end{definicion}

\begin{ejemplo}
Calcular
\[
D=\begin{vmatrix}1&2&0\\0&3&1\\2&0&1\end{vmatrix}.
\]
Por la fórmula anterior:
\begin{align*}
D&=(1)(3)(1)+(2)(1)(2)+(0)(0)(0)\\
&\quad -(0)(3)(2)-(2)(0)(1)-(1)(1)(0)\\
&=3+4=\boxed{7}.
\end{align*}
\end{ejemplo}

%=======================================================
\subsection*{Definición general: menores y cofactores}
%=======================================================
\begin{definicion}[Menor y cofactor]
Sea $A=(a_{ij})$ una matriz cuadrada de orden $n\ge2$. El \emph{menor} $M_{ij}$ de la entrada $a_{ij}$ es el determinante de la matriz de orden $n-1$ obtenida al suprimir la fila $i$ y la columna $j$.

El \emph{cofactor} correspondiente es
\[
C_{ij}=(-1)^{i+j}M_{ij}.
\]
\end{definicion}

Los signos de los cofactores alternan siguiendo el patrón
\[
\begin{pmatrix}
+&-&+&\cdots\\-&+&-&\cdots\\+&-&+&\cdots\\\vdots&\vdots&\vdots&\ddots
\end{pmatrix}.
\]

\begin{teorema}[Expansión por cofactores o de Laplace]
Sea $A=(a_{ij})$ una matriz cuadrada de orden $n\ge2$. Fijando una fila $i$, su determinante puede calcularse como
\[
\boxed{\det(A)=\sum_{j=1}^{n}a_{ij}C_{ij}.}
\]
También puede expandirse a lo largo de cualquier columna $j$:
\[
\det(A)=\sum_{i=1}^{n}a_{ij}C_{ij}.
\]
Para cerrar la definición recursiva se toma $\det([a])=a$ para matrices de orden uno.
\end{teorema}

\begin{observacion}
Para calcular un determinante mediante cofactores conviene elegir una fila o una columna con muchos ceros: se reduce así el número de determinantes menores que hay que evaluar.
\end{observacion}

\begin{ejemplo}[Expansión por la primera fila]
Consideremos nuevamente
\[
A=\begin{pmatrix}1&2&0\\0&3&1\\2&0&1\end{pmatrix}.
\]
Los menores de la primera fila son
\[
M_{11}=\begin{vmatrix}3&1\\0&1\end{vmatrix}=3,\quad
M_{12}=\begin{vmatrix}0&1\\2&1\end{vmatrix}=-2,\quad
M_{13}=\begin{vmatrix}0&3\\2&0\end{vmatrix}=-6.
\]
Los cofactores son $C_{11}=3$, $C_{12}=2$ y $C_{13}=-6$. Por tanto,
\[
\det(A)=1(3)+2(2)+0(-6)=\boxed{7}.
\]
\end{ejemplo}
%=======================================================
\subsection*{Propiedades de los determinantes}
%=======================================================
\begin{proposicion}
Para matrices cuadradas de orden $n$:
\begin{enumerate}
\item $\det(I_n)=1$.
\item $\det(A^T)=\det(A)$.
\item Si $A$ es triangular, $\det(A)=\prod_{i=1}^{n}a_{ii}$.
\item Si dos filas de $A$ son iguales o proporcionales, entonces $\det(A)=0$.
\item Intercambiar dos filas cambia el signo del determinante.
\item Multiplicar una fila por $\lambda$ multiplica el determinante por $\lambda$.
\item Sumar a una fila un múltiplo de otra fila \emph{no cambia} el determinante.
\item $\det(AB)=\det(A)\det(B)$.
\end{enumerate}
\end{proposicion}

\begin{observacion}
Las propiedades anteriores pueden demostrarse rigurosamente a partir de la definición del determinante mediante permutaciones. Aquí las empleamos como resultados del curso; su justificación completa mediante esa construcción combinatoria sería una ampliación teórica adicional.
\end{observacion}

\begin{ejemplo}
Calcular
\[
\det\begin{pmatrix}2&1&0\\4&3&1\\0&2&3\end{pmatrix}.
\]
Aplicamos operaciones que no cambian el determinante:
\[
\begin{pmatrix}2&1&0\\4&3&1\\0&2&3\end{pmatrix}
\xrightarrow{R_2\leftarrow R_2-2R_1}
\begin{pmatrix}2&1&0\\0&1&1\\0&2&3\end{pmatrix}
\xrightarrow{R_3\leftarrow R_3-2R_2}
\begin{pmatrix}2&1&0\\0&1&1\\0&0&1\end{pmatrix}.
\]
La última matriz es triangular superior. Su determinante es el producto de las entradas diagonales:
\[
\boxed{\det(A)=2\cdot1\cdot1=2}.
\]
\end{ejemplo}
%=======================================================
\subsection*{Determinantes e invertibilidad}
%=======================================================
\begin{teorema}[Criterio de invertibilidad]
Para una matriz cuadrada $A$ de orden $n$, las siguientes condiciones son equivalentes:
\begin{enumerate}
\item $A$ es invertible.
\item $\det(A)\ne0$.
\item Las filas de $A$ son linealmente independientes.
\item $\rang(A)=n$.
\item El sistema homogéneo $A\vect{x}=\vect{0}$ sólo tiene la solución trivial.
\end{enumerate}
\end{teorema}

%===============================================================
\section{Matriz adjunta e inversa mediante cofactores}
%===============================================================
\begin{definicion}
Sea $A=(a_{ij})$ cuadrada. La \emph{matriz de cofactores} es $C=(C_{ij})$. Su transpuesta recibe el nombre de \emph{matriz adjunta}:
\[
\boxed{\adj(A)=C^T.}
\]
\end{definicion}

\begin{teorema}
Para cualquier matriz cuadrada $A$,
\[
A\adj(A)=\adj(A)A=\det(A)I_n.
\]
En consecuencia, si $\det(A)\ne0$,
\[
\boxed{A^{-1}=\frac{1}{\det(A)}\adj(A).}
\]
\end{teorema}

\begin{ejemplo}
Para
\[
A=\begin{pmatrix}a&b\\c&d\end{pmatrix},\qquad \det(A)=ad-bc\ne0,
\]
la matriz de cofactores es $\begin{psmallmatrix}d&-c\\-b&a\end{psmallmatrix}$ y su adjunta es $\begin{psmallmatrix}d&-b\\-c&a\end{psmallmatrix}$. Así,
\[
\boxed{A^{-1}=\frac{1}{ad-bc}\begin{pmatrix}d&-b\\-c&a\end{pmatrix}.}
\]
En particular, para $A=\begin{psmallmatrix}2&1\\1&1\end{psmallmatrix}$:
\[
\det(A)=1,\qquad A^{-1}=\begin{pmatrix}1&-1\\-1&2\end{pmatrix},
\]
el mismo resultado obtenido por Gauss--Jordan.
\end{ejemplo}

\begin{ejemplo}[Adjunta e inversa de una matriz $3\times3$]
Consideremos
\[
A=\begin{pmatrix}2&1&0\\1&2&1\\0&1&2\end{pmatrix}.
\]
Expandiendo por la primera fila:
\[
\det(A)=2\begin{vmatrix}2&1\\1&2\end{vmatrix}
-1\begin{vmatrix}1&1\\0&2\end{vmatrix}
=2(3)-2=4.
\]
Calculamos los nueve cofactores:
\[
C=\begin{pmatrix}
3&-2&1\\
-2&4&-2\\
1&-2&3
\end{pmatrix}.
\]
En este ejemplo $C=C^T$, por lo que
\[
\adj(A)=C^T=\begin{pmatrix}
3&-2&1\\-2&4&-2\\1&-2&3
\end{pmatrix}.
\]
Se concluye
\[
\boxed{A^{-1}=\frac14\begin{pmatrix}
3&-2&1\\-2&4&-2\\1&-2&3
\end{pmatrix}.}
\]
Comprobación, multiplicando la primera fila de $A$ por las columnas de la adjunta:
\[
(2,1,0)\adj(A)=(4,0,0).
\]
Por simetría y multiplicación directa en las otras filas, $A\adj(A)=4I_3$, como afirma el teorema.
\end{ejemplo}

%===============================================================

\begin{observacion}[Diferencia entre adjunta y transpuesta]
La matriz de cofactores se denota $C=(C_{ij})$ y la \emph{adjunta clásica} es $\operatorname{adj}(A)=C^T$. No debe confundirse con la transpuesta $A^T$ ni, para matrices complejas, con la conjugada transpuesta $A^*$. Son objetos distintos.
\end{observacion}
%=======================================================
\subsection*{Decidir si conviene Gauss--Jordan o la fórmula de la adjunta}
%=======================================================
Para órdenes dos o tres, la adjunta puede ofrecer un cálculo manual razonable. En órdenes cuatro y cinco, el número de menores suele hacerla laboriosa. Gauss--Jordan suele resultar más sistemático; si sólo nos interesa decidir invertibilidad, calcular $\det(A)$ por eliminación o contar pivotes puede ser más eficiente.
\begin{ejemplo}[Cómo reconocer rápidamente una matriz no invertible]
Sea
\[
A=\begin{pmatrix}1&0&2\\2&0&4\\0&1&3\end{pmatrix}.
\]
La segunda fila es el doble de la primera. Por alternancia, $\det(A)=0$. Al hacer $R_2\leftarrow R_2-2R_1$ aparece una fila nula. El rango es menor que tres; no existe $A^{-1}$ y el sistema homogéneo tiene soluciones distintas de cero. Las tres conclusiones provienen del mismo hecho.
\end{ejemplo}

%=======================================================
\section{Rango y núcleo de una matriz}
%===============================================================
Estos conceptos sirven como puente hacia los espacios vectoriales y las transformaciones lineales. En esta unidad los introduciremos a partir del cálculo con matrices.

\begin{definicion}[Rango]
El \emph{rango} de una matriz $A$, denotado por $\rang(A)$, es el número de pivotes de una forma escalonada obtenida mediante operaciones elementales por filas. Equivalentemente, es el número de columnas linealmente independientes de $A$.
\end{definicion}

\begin{definicion}[Núcleo]
Para $A\in\Mat{m}{n}(\R)$, su \emph{núcleo} es el conjunto de soluciones del sistema homogéneo asociado:
\[
\boxed{\Ker(A)=\{\vect{x}\in\R^n:A\vect{x}=\vect{0}\}.}
\]
\end{definicion}

\begin{observacion}
Las operaciones elementales por filas preservan las soluciones de $A\vect{x}=\vect0$. Por tanto, podemos calcular el núcleo a partir de la forma escalonada reducida. Sin embargo, esas operaciones no preservan, en general, el \emph{espacio de columnas} como subconjunto de $\R^m$; para construir una base de las columnas utilizaremos las columnas pivote de la \emph{matriz original}.
\end{observacion}

\begin{ejemplo}[Rango y núcleo paso a paso]
Sea
\[
A=\begin{pmatrix}1&2&3\\2&4&6\\0&1&1\end{pmatrix}.
\]
Buscamos la forma escalonada reducida:
\begin{align*}
A&\xrightarrow{R_2\leftarrow R_2-2R_1}
\begin{pmatrix}1&2&3\\0&0&0\\0&1&1\end{pmatrix}
\xrightarrow{R_2\leftrightarrow R_3}
\begin{pmatrix}1&2&3\\0&1&1\\0&0&0\end{pmatrix}\\[1ex]
&\xrightarrow{R_1\leftarrow R_1-2R_2}
\begin{pmatrix}1&0&1\\0&1&1\\0&0&0\end{pmatrix}.
\end{align*}
Hay dos pivotes; por tanto,
\[\boxed{\rang(A)=2.}\]
El sistema $A\vect{x}=\vect{0}$ equivale a
\[
\begin{cases}
x_1+x_3=0,\\x_2+x_3=0.
\end{cases}
\]
Tomando $x_3=t$ como variable libre, obtenemos
\[
x_1=-t,\qquad x_2=-t,\qquad x_3=t.
\]
Luego
\[
\boxed{\Ker(A)=\left\{t\begin{pmatrix}-1\\-1\\1\end{pmatrix}:t\in\R\right\}.}
\]
Las columnas pivote son la primera y la segunda; en la matriz original son
\[
\begin{pmatrix}1\\2\\0\end{pmatrix},\qquad
\begin{pmatrix}2\\4\\1\end{pmatrix}.
\]
La tercera es la suma de las dos primeras. Esta observación da una interpretación concreta de $\rang(A)=2$.
\end{ejemplo}

\begin{proposicion}
Para toda matriz $A\in\Mat{m}{n}(\R)$,
\[
0\le\rang(A)\le\min\{m,n\}.
\]
Si $A$ es cuadrada de orden $n$, entonces es invertible si y sólo si $\rang(A)=n$.
\end{proposicion}

\begin{observacion}
Si el núcleo de una matriz $m\times n$ tiene dimensión $d$, el teorema rango--nulidad establece que $\rang(A)+d=n$. En la Unidad 5 demostraremos este resultado en el lenguaje de las transformaciones lineales.
\end{observacion}

%===============================================================
\subsection*{Espacios columna, espacios fila y el significado de rango}
%=======================================================
Las columnas de $A$ pertenecen a $\R^m$. El conjunto de todas sus combinaciones lineales se llama \emph{espacio columna}:
\[
\operatorname{Col}(A)=\{A\vect x:\vect x\in\R^n\}
=\left\{\sum_{j=1}^{n}x_j\vect a_j:x_j\in\R\right\}.
\]
De forma análoga, las combinaciones de filas forman el espacio fila de $A$. El rango es la dimensión de ambos espacios. Podemos interpretar el rango como el número de pivotes (o columnas pivote) de una matriz escalonada.

\begin{proposicion}[Qué conservan las operaciones elementales]
Si $E$ es invertible y $B=EA$, entonces $\Ker(B)=\Ker(A)$ y $\rang(B)=\rang(A)$. El espacio columna \emph{como conjunto de vectores} no tiene por qué conservarse.
\end{proposicion}
\begin{proof}
$B\vect x=\vect0$ equivale a $EA\vect x=\vect0$; multiplicando por $E^{-1}$ equivale a $A\vect x=\vect0$. Las operaciones elementales no cambian el número de pivotes (ni la cantidad de variables libres), por lo que preservan el rango. Pero $\operatorname{Col}(EA)=E\operatorname{Col}(A)$, un conjunto que puede ser diferente de $\operatorname{Col}(A)$.
\end{proof}
\begin{ejemplo}[Las columnas pivote deben tomarse de la matriz original]
Sea $A=\begin{psmallmatrix}1&2&3\\2&4&6\\0&1&1\end{psmallmatrix}$. Su forma reducida por filas tiene columnas pivote primera y segunda. Por tanto,
\[
\operatorname{Col}(A)=\operatorname{span}\left\{\begin{pmatrix}1\\2\\0\end{pmatrix},\begin{pmatrix}2\\4\\1\end{pmatrix}\right\}.
\]
Sería incorrecto afirmar que las columnas $(1,0,0)^T$ y $(0,1,0)^T$ de la forma reducida original son, por este motivo, una base de $\operatorname{Col}(A)$.
\end{ejemplo}
%=======================================================
\subsection*{Teorema rango--nulidad explicado mediante pivotes}
%=======================================================
\begin{definicion}[Nulidad]
Para $A\in\Mat{m}{n}(\R)$, llamamos $\operatorname{nul}(A)$ al número de parámetros independientes que aparecen en la solución general de $A\vect x=\vect0$. En la Unidad 4 reconoceremos este número como $\dim\Ker(A)$.
\end{definicion}
\begin{teorema}[Rango--nulidad para una matriz]
Para una matriz $A$ con $n$ columnas,
\[
\boxed{\rang(A)+\operatorname{nul}(A)=n.}
\]
\end{teorema}
\begin{proof}
Llevemos $A$ a forma escalonada reducida. Si tiene $r$ pivotes, exactamente $r$ variables son básicas y las restantes $n-r$ son libres. Cada elección de las variables libres determina una única solución de las básicas, por lo que la solución general depende de $n-r$ parámetros independientes. En consecuencia, $\operatorname{nul}(A)=n-r$ y $\rang(A)+\operatorname{nul}(A)=r+(n-r)=n$.
\end{proof}
\begin{ejemplo}
Sea
\[
A=\begin{pmatrix}
1&0&2&0&-1\\0&1&-1&0&3\\0&0&0&0&0\\0&0&0&0&0
\end{pmatrix}.
\]
Ya está en forma escalonada reducida, con pivotes en columnas 1 y 2. Por tanto, $\rang(A)=2$ y las columnas 3, 4 y 5 corresponden a variables libres $x_3=s$, $x_4=t$ y $x_5=u$. Las ecuaciones son
\[
x_1+2s-u=0,\qquad x_2-s+3u=0,
\]
por lo que $x_1=-2s+u$ y $x_2=s-3u$. Resulta
\[
\Ker(A)=\left\{s\begin{pmatrix}-2\\1\\1\\0\\0\end{pmatrix}
+t\begin{pmatrix}0\\0\\0\\1\\0\end{pmatrix}
+u\begin{pmatrix}1\\-3\\0\\0\\1\end{pmatrix}:s,t,u\in\R\right\}.
\]
En particular, $\operatorname{nul}(A)=3$ y $2+3=5$, igual al número de columnas.
\end{ejemplo}

%===============================================================
% =============================================================
\section{Ejemplos varios}
%=======================================================

%=======================================================
\subsection*{Suma y combinación lineal: matrices $3\times4$}
%=======================================================
Sean
\[
A=\begin{pmatrix}1&-2&0&3\\2&1&-1&0\\0&4&2&-2\end{pmatrix},\qquad
B=\begin{pmatrix}0&1&5&-1\\3&0&2&4\\-2&1&0&3\end{pmatrix}.
\]
Determinar $2A-B$.
Ambas matrices tienen $3$ filas y $4$ columnas, por lo que la operación está definida y el resultado tendrá dimensión $3\times4$. Primero:
\[
2A=\begin{pmatrix}2&-4&0&6\\4&2&-2&0\\0&8&4&-4\end{pmatrix}.
\]
Ahora se resta cada entrada de $B$ a la entrada correspondiente de $2A$:
\[
2A-B=\begin{pmatrix}
2-0&-4-1&0-5&6-(-1)\\
4-3&2-0&-2-2&0-4\\
0-(-2)&8-1&4-0&-4-3
\end{pmatrix}
=\boxed{\begin{pmatrix}2&-5&-5&7\\1&2&-4&-4\\2&7&4&-7\end{pmatrix}}.
\]
%=======================================================
\subsection*{Producto $2\times4$ por $4\times3$: todas las entradas}
%=======================================================
\[
A=\begin{pmatrix}1&0&2&-1\\2&1&0&3\end{pmatrix},\qquad
B=\begin{pmatrix}1&2&0\\0&-1&3\\2&0&1\\-1&1&2\end{pmatrix}.
\]
El número de columnas de $A$ es $4$, igual al número de filas de $B$. Por ello $AB$ existe y es de orden $2\times3$. Para cada entrada realizamos el producto escalar de una fila con una columna:
\begin{align*}
c_{11}&=1(1)+0(0)+2(2)+(-1)(-1)=6, &c_{12}&=1(2)+0(-1)+2(0)+(-1)(1)=1,\\
c_{13}&=1(0)+0(3)+2(1)+(-1)(2)=0, &c_{21}&=2(1)+1(0)+0(2)+3(-1)=-1,\\
c_{22}&=2(2)+1(-1)+0(0)+3(1)=6, &c_{23}&=2(0)+1(3)+0(1)+3(2)=9.
\end{align*}
Así,
\[\boxed{AB=\begin{pmatrix}6&1&0\\-1&6&9\end{pmatrix}.}\]
\textbf{Advertencia:} en este caso $BA$ \emph{no existe}, porque las dimensiones internas de $(4\times3)(2\times4)$ no coinciden ($3\ne2$).
%=======================================================
\subsection*{Determinante de orden cuatro: reducción conservando su valor}
%=======================================================
Sea
\[
D=\begin{vmatrix}
1&2&0&0\\
2&5&1&0\\
0&1&3&2\\
0&0&2&3
\end{vmatrix}.
\]
La operación $R_2\leftarrow R_2-2R_1$ no cambia el determinante. Después $R_3\leftarrow R_3-R_2$ (con la segunda fila ya modificada) tampoco lo cambia:
\[
D=\begin{vmatrix}1&2&0&0\\0&1&1&0\\0&0&2&2\\0&0&2&3\end{vmatrix}.
\]
Finalmente $R_4\leftarrow R_4-R_3$, y se obtiene una matriz triangular superior:
\[
D=\begin{vmatrix}1&2&0&0\\0&1&1&0\\0&0&2&2\\0&0&0&1\end{vmatrix}
=1\cdot1\cdot2\cdot1=\boxed{2}.
\]
%=======================================================
\subsection*{Rango de una matriz $5\times4$: algoritmo con filas dependientes}
%=======================================================
Sea
\[
A=\begin{pmatrix}
1&2&0&1\\0&1&1&-1\\1&3&1&0\\0&0&1&2\\1&2&1&3
\end{pmatrix}.
\]
Observamos que $R_3=R_1+R_2$ y $R_5=R_1+R_4$. En consecuencia, hacemos explícita esa dependencia:
\[
\begin{aligned}
R_3&\leftarrow R_3-R_1-R_2,\\
R_5&\leftarrow R_5-R_1-R_4.
\end{aligned}
\quad\Longrightarrow\quad
\begin{pmatrix}
1&2&0&1\\0&1&1&-1\\0&0&0&0\\0&0&1&2\\0&0&0&0
\end{pmatrix}.
\]
Intercambiamos $R_3$ y $R_4$, quedando una forma escalonada con pivotes en las columnas $1,2,3$. Por tanto,
\[\boxed{\rang(A)=3.}\]
\textbf{Interpretación:} hay cuatro columnas y sólo tres pivotes; por ello, el sistema homogéneo $A\vect x=0$ tendrá una variable libre. Nótese que $A$ no es cuadrada: no tiene sentido preguntar por $\det(A)$ ni por $A^{-1}$.
%=======================================================
\subsection*{Núcleo de una matriz $3\times5$ con dos parámetros}
%=======================================================
Consideremos
\[
A=\begin{pmatrix}1&2&0&1&-1\\0&1&1&0&2\\0&0&1&1&-1\end{pmatrix}.
\]
Hay tres pivotes, en las tres primeras columnas, y dos variables libres. Resolveremos $A\vect x=0$ mediante sustitución hacia atrás, \emph{sin omitir variables}.
De la tercera fila: $x_3+x_4-x_5=0$, así que $x_3=-x_4+x_5$. De la segunda: $x_2+x_3+2x_5=0$, luego $x_2=x_4-3x_5$. De la primera:
\[
x_1+2(x_4-3x_5)+x_4-x_5=0\quad\Longrightarrow\quad x_1=-3x_4+7x_5.
\]
Tomando $x_4=s$ y $x_5=t$, la solución general es
\[
\vect x=\begin{pmatrix}-3s+7t\\s-3t\\-s+t\\s\\t\end{pmatrix}
=s\begin{pmatrix}-3\\1\\-1\\1\\0\end{pmatrix}
+t\begin{pmatrix}7\\-3\\1\\0\\1\end{pmatrix},\qquad s,t\in\R.
\]
Por tanto,
\[
\boxed{\Ker(A)=\operatorname{span}\left\{
\begin{pmatrix}-3\\1\\-1\\1\\0\end{pmatrix},
\begin{pmatrix}7\\-3\\1\\0\\1\end{pmatrix}\right\}.}
\]
Para comprobar que ambos vectores pertenecen al núcleo, se multiplica $A$ por cada uno y se verifica que se obtiene el vector cero. Hay dos parámetros libres: $\nul(A)=2$ y $\rang(A)+\nul(A)=3+2=5$.

%=======================================================
\section{Banco de ejercicios}
%===============================================================
Los ejercicios están ordenados por tema. Se recomienda resolverlos primero sin software y, después, utilizar R, Python o una calculadora como instrumento de verificación. El apéndice contiene soluciones detalladas de una selección representativa; el banco ampliado incorpora además respuestas para autocorrección.

\begin{Ejer} \textbf{Definición, clasificación e igualdad}
\begin{enumerate}
\item Determine el orden de $A=\begin{psmallmatrix}1&0&2\\-3&4&5\end{psmallmatrix}$ e identifique $a_{12}$, $a_{23}$.
\item Clasifique $A=\begin{psmallmatrix}3&0&0\\0&3&0\\0&0&3\end{psmallmatrix}$; indique todas las categorías de la Unidad 3 a las que pertenece.
\item Determine $x,y,z$ si $\begin{psmallmatrix}x-2&3y\\z+1&7\end{psmallmatrix}=\begin{psmallmatrix}5&-6\\4&7\end{psmallmatrix}$.
\item Dé un ejemplo de matriz simétrica $3\times3$ que no sea diagonal y de matriz triangular inferior no diagonal.
\item Decida cuáles de las matrices $\begin{psmallmatrix}1&2\\2&1\end{psmallmatrix}$, $\begin{psmallmatrix}1&0\\3&1\end{psmallmatrix}$ e $I_2$ son simétricas.
\end{enumerate}
\end{Ejer}

\begin{Ejer} \textbf{Operaciones y transposición}
\begin{enumerate}[start=6]
\item Sean $A=\begin{psmallmatrix}1&2\\3&4\end{psmallmatrix}$ y $B=\begin{psmallmatrix}-1&0\\2&5\end{psmallmatrix}$. Calcule $A+B$, $2A-B$ y $A^T$.
\item Calcule $AB$ y $BA$ para $A=\begin{psmallmatrix}1&2\\0&1\end{psmallmatrix}$ y $B=\begin{psmallmatrix}2&0\\1&3\end{psmallmatrix}$.
\item Si $A$ es $2\times3$, $B$ es $3\times4$ y $C$ es $4\times2$, indique el orden de $AB$, $BC$, $(AB)C$ y $C(AB)$.
\item Calcule $AB$ para $A=\begin{psmallmatrix}2&-1&0\\1&3&2\end{psmallmatrix}$ y $B=\begin{psmallmatrix}1&2\\0&-1\\4&1\end{psmallmatrix}$.
\item Verifique con matrices $2\times2$ concretas la identidad $(AB)^T=B^TA^T$.
\item Demuestre que $(A+B)^T=A^T+B^T$ para matrices de igual orden.
\end{enumerate}
\end{Ejer}

\begin{Ejer} \textbf{Inversas}
\begin{enumerate}[start=12]
\item Mediante Gauss--Jordan, obtenga $A^{-1}$ para $A=\begin{psmallmatrix}3&1\\2&1\end{psmallmatrix}$ y compruebe el resultado.
\item Determine si $A=\begin{psmallmatrix}1&2\\3&6\end{psmallmatrix}$ tiene inversa; justifique mediante eliminación.
\item Calcule la inversa, si existe, de $A=\begin{psmallmatrix}1&0&0\\0&2&0\\0&0&-3\end{psmallmatrix}$.
\item Demuestre que $(A^{-1})^{-1}=A$ cuando $A$ es invertible.
\item Si $A$ y $B$ son invertibles, demuestre que $(AB)^{-1}=B^{-1}A^{-1}$.
\end{enumerate}
\end{Ejer}

\begin{Ejer}
\textbf{Determinantes y cofactores}
\begin{enumerate}[start=17]
\item Calcule $\begin{vmatrix}4&-2\\3&5\end{vmatrix}$.
\item Calcule $\begin{vmatrix}1&2&0\\0&3&1\\2&0&1\end{vmatrix}$ mediante Sarrus y expansión por cofactores.
\item Calcule el determinante de $\begin{psmallmatrix}2&1&0\\4&3&1\\0&2&3\end{psmallmatrix}$ mediante operaciones elementales.
\item Obtenga los menores y cofactores de la primera fila de $A=\begin{psmallmatrix}1&0&2\\3&1&-1\\0&2&4\end{psmallmatrix}$.
\item Sin expandir completamente, determine $\det\begin{psmallmatrix}1&2&3\\2&4&6\\0&1&1\end{psmallmatrix}$ y justifique.
\item Si $\det(A)=3$ para una matriz $3\times3$, determine $\det(2A)$ y $\det(A^T)$.
\item Calcule $\det(A)$ para $A=\begin{psmallmatrix}1&2&0&0\\0&3&1&0\\0&0&-2&5\\0&0&0&4\end{psmallmatrix}$.
\end{enumerate}
\end{Ejer}

\begin{Ejer}
\textbf{Adjunta, rango y núcleo}
\begin{enumerate}[start=24]
\item Mediante la fórmula de la adjunta, obtenga la inversa de $A=\begin{psmallmatrix}3&1\\2&1\end{psmallmatrix}$.
\item Determine $\adj(A)$ y $A^{-1}$, si existe, para $A=\begin{psmallmatrix}2&1&0\\1&2&1\\0&1&2\end{psmallmatrix}$.
\item Determine el rango y el núcleo de $A=\begin{psmallmatrix}1&2&3\\2&4&6\\0&1&1\end{psmallmatrix}$.
\item Obtenga el rango y el núcleo de $A=\begin{psmallmatrix}1&2&1\\0&1&1\end{psmallmatrix}$.
\item Compruebe que si $A$ es una matriz invertible de orden $n$, entonces $\Ker(A)=\{\vect0\}$.
\item Sea $A=\begin{psmallmatrix}1&1\\1&1\end{psmallmatrix}$ y $\vect b=\begin{psmallmatrix}2\\2\end{psmallmatrix}$. Describa todas las soluciones de $A\vect{x}=\vect b$ e indique por qué no puede usarse $A^{-1}\vect b$.
\end{enumerate}
\end{Ejer}

\begin{Ejer}
\textbf{Ejercicios integradores tipo examen}
\begin{enumerate}[start=30]
\item Sean $A=\begin{psmallmatrix}1&2\\-1&0\end{psmallmatrix}$ y $B=\begin{psmallmatrix}2&1\\0&-1\end{psmallmatrix}$. Calcule $AB$, $BA$, $\det(A)$ y determine si $A$ es invertible.
\item Resuelva $A\vect x=\vect b$ para $A=\begin{psmallmatrix}2&1\\1&1\end{psmallmatrix}$, $\vect b=\begin{psmallmatrix}7\\4\end{psmallmatrix}$ por (a) eliminación y (b) matriz inversa.
\item Sea $A=\begin{psmallmatrix}1&a\\2&4\end{psmallmatrix}$. Determine los valores de $a$ para los que $A$ no es invertible.
\item Determine $\rang(A)$ y $\Ker(A)$ para $A=\begin{psmallmatrix}1&1&0\\0&1&1\\1&2&1\end{psmallmatrix}$.
\item Calcule el determinante $\begin{vmatrix}1&0&2&0\\0&2&0&1\\0&0&-1&3\\0&0&0&4\end{vmatrix}$ y determine si la matriz correspondiente es invertible.
\end{enumerate}
\end{Ejer}


\begin{Ejer}
\textbf{Indique orden, número de entradas y valores de las entradas señaladas}
\begin{enumerate}[start=35]

\item Para $A=\begin{pmatrix}0&0&2&-3\end{pmatrix}$, determine su orden, el número de entradas, $a_{12}$ y $a_{13}$.

\item Para $A=\begin{pmatrix}-2&-2&1\\-3&2&-2\end{pmatrix}$, determine su orden, el número de entradas, $a_{12}$ y $a_{11}$.

\item Para $A=\begin{pmatrix}-2&1\\-2&-3\\3&-3\end{pmatrix}$, determine su orden, el número de entradas, $a_{12}$ y $a_{12}$.

\item Para $A=\begin{pmatrix}-3&-3&3&0\\-3&-1&2&2\\2&-3&2&2\end{pmatrix}$, determine su orden, el número de entradas, $a_{11}$ y $a_{33}$.

\item Para $A=\begin{pmatrix}0&-1&2\\0&3&2\\-1&3&0\\-3&1&3\end{pmatrix}$, determine su orden, el número de entradas, $a_{32}$ y $a_{13}$.

\item Para $A=\begin{pmatrix}2&0\\1&-3\\3&-2\\3&0\\0&-1\end{pmatrix}$, determine su orden, el número de entradas, $a_{31}$ y $a_{12}$.

\item Para $A=\begin{pmatrix}-2&-3&3&-2&1\\3&1&-1&0&0\end{pmatrix}$, determine su orden, el número de entradas, $a_{22}$ y $a_{21}$.

\item Para $A=\begin{pmatrix}0&1&-2&-3\\0&2&3&-2\\1&3&2&3\\3&1&1&2\end{pmatrix}$, determine su orden, el número de entradas, $a_{11}$ y $a_{33}$.

\item Para $A=\begin{pmatrix}-3&-3&-2\\-2&-3&-1\\-1&1&0\\-3&0&-3\\-2&-3&-2\end{pmatrix}$, determine su orden, el número de entradas, $a_{53}$ y $a_{22}$.

\item Para $A=\begin{pmatrix}2&-3&-2&-1&1\\-1&0&2&-1&0\\2&-1&0&-2&-3\end{pmatrix}$, determine su orden, el número de entradas, $a_{11}$ y $a_{14}$.

\item Para $A=\begin{pmatrix}-2&-2&-2&2\\-1&1&-2&3\\0&1&-2&-1\\-2&-1&-3&-1\\1&-3&-2&0\end{pmatrix}$, determine su orden, el número de entradas, $a_{23}$ y $a_{52}$.

\item Para $A=\begin{pmatrix}-3&1&-3&1&-1\\-3&-2&2&3&-1\\3&1&3&-3&-3\\0&1&-2&-1&2\\0&0&-1&0&3\end{pmatrix}$, determine su orden, el número de entradas, $a_{11}$ y $a_{31}$.

\end{enumerate}
\end{Ejer}

\begin{Ejer}
\textbf{Iguale las entradas en la misma posición para determinar $x$ e $y$. Verifique sustituyendo al final.}
\begin{enumerate}[start=47]

\item Determine $x,y\in\mathbb R$ si $\begin{pmatrix}-1&y - 1\\-2&x - 2\end{pmatrix}=\begin{pmatrix}-1&2\\-2&-3\end{pmatrix}$.

\item Determine $x,y\in\mathbb R$ si $\begin{pmatrix}3&1&0\\y - 2&x + 1&0\end{pmatrix}=\begin{pmatrix}3&1&0\\1&-2&0\end{pmatrix}$.

\item Determine $x,y\in\mathbb R$ si $\begin{pmatrix}x + 1&3\\-2&2\\-1&y - 1\end{pmatrix}=\begin{pmatrix}-3&3\\-2&2\\-1&1\end{pmatrix}$.

\item Determine $x,y\in\mathbb R$ si $\begin{pmatrix}0&-1&3\\x + 2&1&3\\y - 2&3&-2\end{pmatrix}=\begin{pmatrix}0&-1&3\\1&1&3\\0&3&-2\end{pmatrix}$.

\item Determine $x,y\in\mathbb R$ si $\begin{pmatrix}x + 2&-3&-1&3\\-3&y - 1&-1&-1\end{pmatrix}=\begin{pmatrix}1&-3&-1&3\\-3&-3&-1&-1\end{pmatrix}$.

\item Determine $x,y\in\mathbb R$ si $\begin{pmatrix}-2&2\\-3&0\\1&x + 3\\-3&y - 1\end{pmatrix}=\begin{pmatrix}-2&2\\-3&0\\1&-3\\-3&1\end{pmatrix}$.

\item Determine $x,y\in\mathbb R$ si $\begin{pmatrix}-1&x + 2&3&-2\\-1&-2&-3&0\\3&y + 1&0&0\end{pmatrix}=\begin{pmatrix}-1&-1&3&-2\\-1&-2&-3&0\\3&-3&0&0\end{pmatrix}$.

\item Determine $x,y\in\mathbb R$ si $\begin{pmatrix}-1&-2&-3\\0&1&-1\\-3&2&y - 2\\3&x + 1&3\end{pmatrix}=\begin{pmatrix}-1&-2&-3\\0&1&-1\\-3&2&-1\\3&3&3\end{pmatrix}$.

\item Determine $x,y\in\mathbb R$ si $\begin{pmatrix}-3&-3\\y - 1&-3\\-1&-3\\x - 1&3\\-2&-3\end{pmatrix}=\begin{pmatrix}-3&-3\\-2&-3\\-1&-3\\-3&3\\-2&-3\end{pmatrix}$.

\item Determine $x,y\in\mathbb R$ si $\begin{pmatrix}y + 2&0&2&2&0\\x + 1&3&3&-2&-3\end{pmatrix}=\begin{pmatrix}2&0&2&2&0\\0&3&3&-2&-3\end{pmatrix}$.

\item Determine $x,y\in\mathbb R$ si $\begin{pmatrix}0&0&-2&3\\y - 2&-2&-1&-2\\-1&1&-2&0\\-3&1&x + 1&-2\end{pmatrix}=\begin{pmatrix}0&0&-2&3\\-1&-2&-1&-2\\-1&1&-2&0\\-3&1&-2&-2\end{pmatrix}$.

\item Determine $x,y\in\mathbb R$ si $\begin{pmatrix}3&-3&-3&3&3\\2&1&1&0&3\\-2&2&2&-3&-1\\y + 1&2&-3&2&x + 2\\-1&-3&1&3&-2\end{pmatrix}=\begin{pmatrix}3&-3&-3&3&3\\2&1&1&0&3\\-2&2&2&-3&-1\\-2&2&-3&2&0\\-1&-3&1&3&-2\end{pmatrix}$.

\end{enumerate}
\end{Ejer}

\begin{Ejer}

\textbf{Para cada matriz identifique todas las categorías aplicables: fila, columna, rectangular, cuadrada, diagonal, identidad, triangular, simétrica y nula.}
\begin{enumerate}[start=59]

\item Clasifique $A=\begin{pmatrix}2&0\\0&3\end{pmatrix}$.

\item Clasifique $A=\begin{pmatrix}1&0\\0&1\end{pmatrix}$.

\item Clasifique $A=\begin{pmatrix}1&2&0\\2&1&2\\0&2&1\end{pmatrix}$.

\item Clasifique $A=\begin{pmatrix}1&2&0\\0&1&0\\0&0&1\end{pmatrix}$.

\item Clasifique $A=\begin{pmatrix}1&0&0&0\\-2&1&0&0\\0&0&1&0\\0&0&0&1\end{pmatrix}$.

\item Clasifique $A=\begin{pmatrix}2&0&-3\\-1&-2&-3\\3&2&1\\0&3&3\end{pmatrix}$.

\item Clasifique $A=\begin{pmatrix}0&0&0&0&0\\0&0&0&0&0\\0&0&0&0&0\\0&0&0&0&0\\0&0&0&0&0\end{pmatrix}$.

\item Clasifique $A=\begin{pmatrix}1&1&0&0&0\\2&1&0&0&0\\0&0&1&0&0\\0&0&0&1&0\\0&0&0&0&1\end{pmatrix}$.

\item Clasifique $A=\begin{pmatrix}-1&0\\0&-1\end{pmatrix}$.

\item Clasifique $A=\begin{pmatrix}1&0\\0&1\end{pmatrix}$.

\item Clasifique $A=\begin{pmatrix}1&2&0\\2&1&2\\0&2&1\end{pmatrix}$.

\item Clasifique $A=\begin{pmatrix}1&2&0\\0&1&0\\0&0&1\end{pmatrix}$.

\end{enumerate}
\end{Ejer}

\begin{Ejer}

\textbf{En todos los ejercicios calcule $2A-B$, trabajando entrada por entrada. Esta repetición busca automatizar el cálculo sin cambiar la técnica.}
\begin{enumerate}[start=71]

\item Si $A=\begin{pmatrix}-2&-1\\2&1\end{pmatrix}$ y $B=\begin{pmatrix}-1&-2\\-2&1\end{pmatrix}$, calcule $2A-B$.

\item Si $A=\begin{pmatrix}-1&-1&2\\1&-1&0\end{pmatrix}$ y $B=\begin{pmatrix}1&-2&-2\\2&-2&-1\end{pmatrix}$, calcule $2A-B$.

\item Si $A=\begin{pmatrix}1&1\\-2&1\\1&1\end{pmatrix}$ y $B=\begin{pmatrix}1&2\\0&2\\1&-2\end{pmatrix}$, calcule $2A-B$.

\item Si $A=\begin{pmatrix}-2&-2&-2\\-1&-2&1\\0&-1&1\end{pmatrix}$ y $B=\begin{pmatrix}1&1&-1\\2&0&-1\\1&2&0\end{pmatrix}$, calcule $2A-B$.

\item Si $A=\begin{pmatrix}-1&2\\-1&1\\-2&-2\\-2&1\end{pmatrix}$ y $B=\begin{pmatrix}1&1\\1&0\\2&-2\\-1&1\end{pmatrix}$, calcule $2A-B$.

\item Si $A=\begin{pmatrix}0&-2&-2&-2\\-2&0&0&2\end{pmatrix}$ y $B=\begin{pmatrix}-1&-1&-1&-1\\-1&2&-1&-2\end{pmatrix}$, calcule $2A-B$.

\item Si $A=\begin{pmatrix}1&1&2&2\\0&-1&0&2\\-1&1&-1&0\end{pmatrix}$ y $B=\begin{pmatrix}-2&2&0&1\\0&2&2&1\\1&0&-1&-1\end{pmatrix}$, calcule $2A-B$.

\item Si $A=\begin{pmatrix}1&-2&2\\2&2&2\\-2&-2&0\\-2&2&0\end{pmatrix}$ y $B=\begin{pmatrix}-1&-1&-1\\0&0&2\\1&-2&-1\\0&-2&-2\end{pmatrix}$, calcule $2A-B$.

\item Si $A=\begin{pmatrix}-1&1\\-2&-2\\1&2\\0&2\\2&-1\end{pmatrix}$ y $B=\begin{pmatrix}-1&2\\2&1\\-2&-1\\-1&-1\\2&-2\end{pmatrix}$, calcule $2A-B$.

\item Si $A=\begin{pmatrix}0&0&-1&-2&-1\\0&0&0&1&-1\end{pmatrix}$ y $B=\begin{pmatrix}1&2&1&-1&2\\1&-1&-1&1&2\end{pmatrix}$, calcule $2A-B$.

\item Si $A=\begin{pmatrix}1&0&1&-2\\-2&2&1&-2\\2&-1&2&2\\-1&0&1&0\end{pmatrix}$ y $B=\begin{pmatrix}1&-2&1&-2\\0&1&0&-1\\0&1&-1&0\\2&0&-2&-1\end{pmatrix}$, calcule $2A-B$.

\item Si $A=\begin{pmatrix}-2&1&1\\1&-1&1\\1&-1&1\\-2&-1&1\\2&1&-1\end{pmatrix}$ y $B=\begin{pmatrix}-2&2&1\\-1&2&0\\-1&1&-1\\-2&-1&2\\-1&1&1\end{pmatrix}$, calcule $2A-B$.

\item Si $A=\begin{pmatrix}-1&-2&-2&-1&0\\-2&-1&0&-1&-1\\-2&-2&-2&-1&2\end{pmatrix}$ y $B=\begin{pmatrix}-1&0&2&0&-2\\2&-1&-1&0&2\\0&0&2&-2&1\end{pmatrix}$, calcule $2A-B$.

\item Si $A=\begin{pmatrix}1&2&2&0&2\\-2&-1&2&-2&-2\\1&-2&2&1&1\\-1&-2&2&1&2\\-2&0&-1&0&0\end{pmatrix}$ y $B=\begin{pmatrix}-2&-2&1&1&-2\\-2&-1&1&2&0\\2&1&0&2&-1\\-1&2&-1&1&0\\1&0&2&1&-1\end{pmatrix}$, calcule $2A-B$.

\end{enumerate}
\end{Ejer}

\begin{Ejer}

\textbf{Calcule la transpuesta y señale su orden; después indique si la suma $A+A^T$ está definida.}
\begin{enumerate}[start=85]

\item Para $A=\begin{pmatrix}-1&2&0\\-1&-1&1\end{pmatrix}$, calcule $A^T$ y decida si existe $A+A^T$.

\item Para $A=\begin{pmatrix}-1&2\\2&-2\\0&0\end{pmatrix}$, calcule $A^T$ y decida si existe $A+A^T$.

\item Para $A=\begin{pmatrix}-1&-1&-1\\-2&0&-2\\-2&-1&-2\end{pmatrix}$, calcule $A^T$ y decida si existe $A+A^T$.

\item Para $A=\begin{pmatrix}1&0\\0&1\\-2&0\\2&1\end{pmatrix}$, calcule $A^T$ y decida si existe $A+A^T$.

\item Para $A=\begin{pmatrix}1&2&-1&-1\\-1&0&-1&-2\end{pmatrix}$, calcule $A^T$ y decida si existe $A+A^T$.

\item Para $A=\begin{pmatrix}1&2&2&-2\\2&-1&-1&0\\2&-2&0&-2\\-2&-1&-1&-2\end{pmatrix}$, calcule $A^T$ y decida si existe $A+A^T$.

\item Para $A=\begin{pmatrix}0&1\\2&1\\-2&0\\1&1\\-2&-2\end{pmatrix}$, calcule $A^T$ y decida si existe $A+A^T$.

\item Para $A=\begin{pmatrix}1&0&-1&2&-1\\0&-2&1&2&-1\end{pmatrix}$, calcule $A^T$ y decida si existe $A+A^T$.

\item Para $A=\begin{pmatrix}-2&2&0&-1\\0&0&2&1\\-1&1&1&0\\2&1&2&-1\\1&-2&-1&0\end{pmatrix}$, calcule $A^T$ y decida si existe $A+A^T$.

\item Para $A=\begin{pmatrix}1&1&1&-2&-1\\2&-2&-1&0&2\\-2&-1&-2&1&2\\-1&1&-1&-1&2\end{pmatrix}$, calcule $A^T$ y decida si existe $A+A^T$.

\item Para $A=\begin{pmatrix}-1&1&-1&1&0\\2&-1&0&-1&-2\\1&2&1&-1&2\\-1&-1&1&1&-2\\0&-1&-2&1&2\end{pmatrix}$, calcule $A^T$ y decida si existe $A+A^T$.

\item Para $A=\begin{pmatrix}1&-1&2&-2&-1\\-2&0&0&2&-2\\0&-1&-1&0&-1\end{pmatrix}$, calcule $A^T$ y decida si existe $A+A^T$.

\end{enumerate}
\end{Ejer}

\begin{Ejer}

\textbf{Sin realizar productos numéricos, determine cuáles existen y el orden del resultado. Si no existen, explique qué dimensiones internas no coinciden.}
\begin{enumerate}[start=97]

\item Sean $A$ de orden $2\times3$ y $B$ de orden $3\times4$. ¿Existe $AB$? ¿Existe $BA$? Indique los órdenes cuando existan.

\item Sean $A$ de orden $3\times2$ y $B$ de orden $2\times5$. ¿Existe $AB$? ¿Existe $BA$? Indique los órdenes cuando existan.

\item Sean $A$ de orden $4\times3$ y $B$ de orden $4\times2$. ¿Existe $AB$? ¿Existe $BA$? Indique los órdenes cuando existan.

\item Sean $A$ de orden $2\times5$ y $B$ de orden $5\times3$. ¿Existe $AB$? ¿Existe $BA$? Indique los órdenes cuando existan.

\item Sean $A$ de orden $5\times4$ y $B$ de orden $4\times2$. ¿Existe $AB$? ¿Existe $BA$? Indique los órdenes cuando existan.

\item Sean $A$ de orden $3\times5$ y $B$ de orden $4\times3$. ¿Existe $AB$? ¿Existe $BA$? Indique los órdenes cuando existan.

\item Sean $A$ de orden $4\times2$ y $B$ de orden $2\times3$. ¿Existe $AB$? ¿Existe $BA$? Indique los órdenes cuando existan.

\item Sean $A$ de orden $3\times3$ y $B$ de orden $3\times4$. ¿Existe $AB$? ¿Existe $BA$? Indique los órdenes cuando existan.

\item Sean $A$ de orden $2\times4$ y $B$ de orden $2\times4$. ¿Existe $AB$? ¿Existe $BA$? Indique los órdenes cuando existan.

\item Sean $A$ de orden $5\times5$ y $B$ de orden $5\times2$. ¿Existe $AB$? ¿Existe $BA$? Indique los órdenes cuando existan.

\item Sean $A$ de orden $1\times5$ y $B$ de orden $5\times1$. ¿Existe $AB$? ¿Existe $BA$? Indique los órdenes cuando existan.

\item Sean $A$ de orden $5\times1$ y $B$ de orden $1\times5$. ¿Existe $AB$? ¿Existe $BA$? Indique los órdenes cuando existan.

\item Sean $A$ de orden $4\times5$ y $B$ de orden $5\times4$. ¿Existe $AB$? ¿Existe $BA$? Indique los órdenes cuando existan.

\item Sean $A$ de orden $5\times3$ y $B$ de orden $3\times5$. ¿Existe $AB$? ¿Existe $BA$? Indique los órdenes cuando existan.

\item Sean $A$ de orden $3\times4$ y $B$ de orden $5\times2$. ¿Existe $AB$? ¿Existe $BA$? Indique los órdenes cuando existan.

\item Sean $A$ de orden $2\times3$ y $B$ de orden $2\times5$. ¿Existe $AB$? ¿Existe $BA$? Indique los órdenes cuando existan.

\end{enumerate}
\end{Ejer}

\begin{Ejer}

\textbf{Calcule $AB$, mostrando los productos escalares que conducen a sus entradas. Varíe sólo el tamaño: la técnica siempre es la misma.}
\begin{enumerate}[start=113]

\item Calcule $AB$ para $A=\begin{pmatrix}1&0\\0&2\end{pmatrix}$ y $B=\begin{pmatrix}0&-1\\-2&-2\end{pmatrix}$.

\item Calcule $AB$ para $A=\begin{pmatrix}1&2&0\\0&0&2\end{pmatrix}$ y $B=\begin{pmatrix}-2&0\\0&-1\\0&-1\end{pmatrix}$.

\item Calcule $AB$ para $A=\begin{pmatrix}1&-2\\0&-2\\0&-1\end{pmatrix}$ y $B=\begin{pmatrix}-1&0&-2&-1\\-1&2&2&2\end{pmatrix}$.

\item Calcule $AB$ para $A=\begin{pmatrix}2&2&-1&-1\\1&2&2&2\end{pmatrix}$ y $B=\begin{pmatrix}-2&1&1\\-1&-1&1\\-2&1&2\\0&1&2\end{pmatrix}$.

\item Calcule $AB$ para $A=\begin{pmatrix}1&1\\2&-1\\2&2\\-2&-1\end{pmatrix}$ y $B=\begin{pmatrix}2&-2&-2\\1&-2&2\end{pmatrix}$.

\item Calcule $AB$ para $A=\begin{pmatrix}-1&0&-1\\1&-1&-1\\-1&-2&1\end{pmatrix}$ y $B=\begin{pmatrix}0&0&-1\\-1&-2&2\\-1&-1&1\end{pmatrix}$.

\item Calcule $AB$ para $A=\begin{pmatrix}2&2&0\\-1&0&1\\-2&2&0\\-1&1&-2\end{pmatrix}$ y $B=\begin{pmatrix}2&-2\\-1&2\\0&2\end{pmatrix}$.

\item Calcule $AB$ para $A=\begin{pmatrix}0&1\\-2&-1\\1&-2\\-2&-1\\0&1\end{pmatrix}$ y $B=\begin{pmatrix}-1&-2&-2\\2&-2&0\end{pmatrix}$.

\item Calcule $AB$ para $A=\begin{pmatrix}1&-2&0&0&1\\-1&-1&-2&0&-2\\2&0&-2&2&1\end{pmatrix}$ y $B=\begin{pmatrix}2&0\\0&-1\\0&2\\2&1\\-2&0\end{pmatrix}$.

\item Calcule $AB$ para $A=\begin{pmatrix}2&-2&-2&-2\\0&2&0&-1\\2&2&1&1\\0&0&-1&-1\end{pmatrix}$ y $B=\begin{pmatrix}-2&-2\\2&1\\0&2\\-2&2\end{pmatrix}$.

\item Calcule $AB$ para $A=\begin{pmatrix}2&2&2\\1&-2&2\\0&2&-1\\2&-2&-1\\2&1&2\end{pmatrix}$ y $B=\begin{pmatrix}1&-1&1&0\\0&-2&-1&-2\\-1&1&-2&1\end{pmatrix}$.

\item Calcule $AB$ para $A=\begin{pmatrix}-2&0&0&-2&2\\-2&-1&-1&-1&0\end{pmatrix}$ y $B=\begin{pmatrix}-1&2&2&1&2\\2&-1&-2&2&2\\-1&-1&1&2&-1\\1&-2&-2&-2&2\\2&1&2&0&2\end{pmatrix}$.

\item Calcule $AB$ para $A=\begin{pmatrix}-1&2&0&-2&-1\\0&2&2&-1&2\\-2&-2&-2&0&-2\\1&0&-2&0&0\\-1&-2&-1&1&2\end{pmatrix}$ y $B=\begin{pmatrix}0&1\\-1&-1\\-1&-2\\-2&-2\\0&2\end{pmatrix}$.

\item Calcule $AB$ para $A=\begin{pmatrix}-1&-2&-1&-2&0\\2&0&-2&-2&-1\\2&-1&1&0&2\\-2&0&2&1&-2\end{pmatrix}$ y $B=\begin{pmatrix}0&0&-2&2\\-1&-2&-2&1\\-1&2&-1&-1\\2&-2&-2&2\\-2&-2&-1&0\end{pmatrix}$.
\end{enumerate}
\end{Ejer}

\begin{Ejer}

\textbf{Forme $[A\mid I_n]$, registre cada operación elemental y verifique $AA^{-1}=I_n$. Los ejemplos de orden $4$ y $5$ se eligieron para mantener operaciones manejables.}\par\medskip
\begin{enumerate}[start=127]

\item Encuentre $A^{-1}$ mediante Gauss--Jordan, para $A=\begin{pmatrix}1&0\\2&1\end{pmatrix}$.

\item Encuentre $A^{-1}$ mediante Gauss--Jordan, para $A=\begin{pmatrix}5&2\\2&1\end{pmatrix}$.

\item Encuentre $A^{-1}$ mediante Gauss--Jordan, para $A=\begin{pmatrix}-3&2\\-2&1\end{pmatrix}$.

\item Encuentre $A^{-1}$ mediante Gauss--Jordan, para $A=\begin{pmatrix}1&0\\0&1\end{pmatrix}$.

\item Encuentre $A^{-1}$ mediante Gauss--Jordan, para $A=\begin{pmatrix}-3&2&0\\-2&1&0\\2&0&1\end{pmatrix}$.

\item Encuentre $A^{-1}$ mediante Gauss--Jordan, para $A=\begin{pmatrix}1&0&2\\0&1&-2\\0&-2&5\end{pmatrix}$.

\item Encuentre $A^{-1}$ mediante Gauss--Jordan, para $A=\begin{pmatrix}1&0&0\\0&1&-1\\0&0&1\end{pmatrix}$.

\item Encuentre $A^{-1}$ mediante Gauss--Jordan, para $A=\begin{pmatrix}1&1&0\\-1&0&0\\0&2&1\end{pmatrix}$.

\item Encuentre $A^{-1}$ mediante Gauss--Jordan, para $A=\begin{pmatrix}1&2&0&0\\0&1&0&0\\0&0&1&-1\\0&0&0&1\end{pmatrix}$.

\item Encuentre $A^{-1}$ mediante Gauss--Jordan, para $A=\begin{pmatrix}1&0&-2&0\\-1&1&0&0\\0&0&1&0\\0&0&0&1\end{pmatrix}$.

\item Encuentre $A^{-1}$ mediante Gauss--Jordan, para $A=\begin{pmatrix}1&0&0&1\\0&1&0&0\\0&0&1&0\\0&0&-1&1\end{pmatrix}$.

\item Encuentre $A^{-1}$ mediante Gauss--Jordan, para $A=\begin{pmatrix}1&0&0&0&0\\0&3&0&2&0\\0&0&1&0&0\\0&1&0&1&0\\0&0&0&0&1\end{pmatrix}$.

\end{enumerate}
\end{Ejer}

\begin{Ejer}

\textbf{En todos los casos aplique la misma identidad $ad-bc$, con especial atención a los signos.}
\begin{enumerate}[start=139]

\item Calcule $\det(A)$ para $A=\begin{pmatrix}-4&-4\\4&-4\end{pmatrix}$.

\item Calcule $\det(A)$ para $A=\begin{pmatrix}4&-2\\4&4\end{pmatrix}$.

\item Calcule $\det(A)$ para $A=\begin{pmatrix}3&2\\-1&-2\end{pmatrix}$.

\item Calcule $\det(A)$ para $A=\begin{pmatrix}-3&-4\\3&4\end{pmatrix}$.

\item Calcule $\det(A)$ para $A=\begin{pmatrix}2&-1\\0&2\end{pmatrix}$.

\item Calcule $\det(A)$ para $A=\begin{pmatrix}-1&3\\2&-1\end{pmatrix}$.

\item Calcule $\det(A)$ para $A=\begin{pmatrix}1&-1\\-3&-2\end{pmatrix}$.

\item Calcule $\det(A)$ para $A=\begin{pmatrix}3&3\\2&-1\end{pmatrix}$.

\item Calcule $\det(A)$ para $A=\begin{pmatrix}-1&-3\\-1&2\end{pmatrix}$.

\item Calcule $\det(A)$ para $A=\begin{pmatrix}1&1\\-1&1\end{pmatrix}$.

\item Calcule $\det(A)$ para $A=\begin{pmatrix}2&4\\3&-2\end{pmatrix}$.

\item Calcule $\det(A)$ para $A=\begin{pmatrix}4&2\\-4&1\end{pmatrix}$.

\end{enumerate}
\end{Ejer}

\begin{Ejer}

\textbf{Calcule cada determinante }
\begin{enumerate}[start=151]

\item Calcule $\det(A)$ para $A=\begin{pmatrix}2&-2&0\\3&3&0\\2&1&3\end{pmatrix}$.

\item Calcule $\det(A)$ para $A=\begin{pmatrix}2&-2&-3\\2&-1&-3\\1&-2&-1\end{pmatrix}$.

\item Calcule $\det(A)$ para $A=\begin{pmatrix}3&3&-1\\3&0&1\\-2&-1&-2\end{pmatrix}$.

\item Calcule $\det(A)$ para $A=\begin{pmatrix}1&-3&2\\-3&-2&2\\1&-1&-1\end{pmatrix}$.

\item Calcule $\det(A)$ para $A=\begin{pmatrix}1&-1&2\\0&-2&1\\1&2&3\end{pmatrix}$.

\item Calcule $\det(A)$ para $A=\begin{pmatrix}-1&1&1\\-2&-1&1\\2&0&1\end{pmatrix}$.

\item Calcule $\det(A)$ para $A=\begin{pmatrix}-1&3&0\\0&-1&2\\2&-1&-2\end{pmatrix}$.

\item Calcule $\det(A)$ para $A=\begin{pmatrix}-3&-3&1\\3&-2&-3\\0&0&-3\end{pmatrix}$.

\item Calcule $\det(A)$ para $A=\begin{pmatrix}2&2&-3\\0&2&3\\0&3&-1\end{pmatrix}$.

\item Calcule $\det(A)$ para $A=\begin{pmatrix}-2&3&-1\\-2&0&0\\0&-1&0\end{pmatrix}$.

\item Calcule $\det(A)$ para $A=\begin{pmatrix}-3&-2&1\\2&2&-3\\1&-3&0\end{pmatrix}$.

\item Calcule $\det(A)$ para $A=\begin{pmatrix}3&3&3\\-1&2&1\\2&-2&2\end{pmatrix}$.

\end{enumerate}
\end{Ejer}

\begin{Ejer}

\textbf{Antes de expandir, busque ceros o utilice eliminación por filas, registrando cambios de signo y factores escalares.}
\begin{enumerate}[start=163]

\item Calcule $\det(A)$ para $A=\begin{pmatrix}2&-2&-4&0\\0&-1&0&0\\0&0&-2&-1\\0&0&0&3\end{pmatrix}$.

\item Calcule $\det(A)$ para $A=\begin{pmatrix}1&0&1&-2\\0&1&0&0\\0&0&-2&2\\0&0&0&1\end{pmatrix}$.

\item Calcule $\det(A)$ para $A=\begin{pmatrix}0&0&1&0\\0&-1&0&0\\-2&1&0&0\\0&0&0&2\end{pmatrix}$.

\item Calcule $\det(A)$ para $A=\begin{pmatrix}0&0&0&1\\0&-1&0&0\\0&0&-2&0\\1&-2&0&0\end{pmatrix}$.

\item Calcule $\det(A)$ para $A=\begin{pmatrix}1&0&0&0\\-1&-1&0&-1\\0&0&1&0\\0&0&0&3\end{pmatrix}$.

\item Calcule $\det(A)$ para $A=\begin{pmatrix}-2&0&2&0\\0&3&0&2\\0&0&1&0\\0&0&0&1\end{pmatrix}$.

\item Calcule $\det(A)$ para $A=\begin{pmatrix}-2&0&2&2&2\\0&2&-2&-2&-2\\0&0&3&0&0\\0&0&0&-1&0\\0&0&0&0&2\end{pmatrix}$.

\item Calcule $\det(A)$ para $A=\begin{pmatrix}-1&0&0&0&2\\0&-2&0&-2&0\\0&0&-2&2&0\\0&0&0&-1&0\\0&0&0&0&1\end{pmatrix}$.

\item Calcule $\det(A)$ para $A=\begin{pmatrix}-1&0&0&0&0\\0&2&0&0&0\\0&0&1&0&0\\0&-2&0&2&0\\0&0&0&0&1\end{pmatrix}$.

\item Calcule $\det(A)$ para $A=\begin{pmatrix}-2&0&0&0&2\\0&3&0&0&-2\\0&0&1&0&-1\\0&0&0&-2&0\\0&-3&0&0&4\end{pmatrix}$.

\item Calcule $\det(A)$ para $A=\begin{pmatrix}-2&0&0&0&1\\0&-1&0&0&1\\0&0&1&2&0\\0&0&0&1&4\\0&0&0&0&3\end{pmatrix}$.

\item Calcule $\det(A)$ para $A=\begin{pmatrix}3&0&0&0&0\\-6&-1&2&0&1\\0&0&1&0&1\\0&0&0&-2&-1\\0&0&0&0&-1\end{pmatrix}$.

\end{enumerate}
\end{Ejer}

\begin{Ejer}

\textbf{Para cada matriz calcule el menor $M_{ij}$ y el cofactor $C_{ij}$ de la entrada indicada. Explique la diferencia entre el menor y el cofactor.}
\begin{enumerate}[start=175]

\item Para $A=\begin{pmatrix}-2&2&2\\1&-2&-2\\-2&2&-1\end{pmatrix}$, encuentre $M_{3,3}$ y $C_{3,3}$.

\item Para $A=\begin{pmatrix}0&0&1\\1&-1&1\\0&-2&2\end{pmatrix}$, encuentre $M_{1,2}$ y $C_{1,2}$.

\item Para $A=\begin{pmatrix}1&-1&-1\\0&-2&-2\\0&1&-1\end{pmatrix}$, encuentre $M_{3,3}$ y $C_{3,3}$.

\item Para $A=\begin{pmatrix}1&1&1\\-1&-1&1\\-1&-1&-2\end{pmatrix}$, encuentre $M_{2,3}$ y $C_{2,3}$.

\item Para $A=\begin{pmatrix}2&0&-2\\1&-1&-1\\-2&0&-1\end{pmatrix}$, encuentre $M_{3,2}$ y $C_{3,2}$.

\item Para $A=\begin{pmatrix}1&2&1\\-1&-2&2\\-1&1&0\end{pmatrix}$, encuentre $M_{2,2}$ y $C_{2,2}$.

\item Para $A=\begin{pmatrix}-2&2&-2\\0&2&2\\2&2&0\end{pmatrix}$, encuentre $M_{3,2}$ y $C_{3,2}$.

\item Para $A=\begin{pmatrix}-1&0&0&0\\1&-2&2&-2\\1&0&2&-2\\2&1&-1&-1\end{pmatrix}$, encuentre $M_{1,4}$ y $C_{1,4}$.

\item Para $A=\begin{pmatrix}1&0&1&-1\\-2&1&-2&-2\\-2&-1&-1&2\\-1&0&2&-2\end{pmatrix}$, encuentre $M_{4,2}$ y $C_{4,2}$.

\item Para $A=\begin{pmatrix}-1&-2&-2&0\\0&-2&0&2\\-1&-1&-1&0\\-2&1&1&-2\end{pmatrix}$, encuentre $M_{4,2}$ y $C_{4,2}$.

\item Para $A=\begin{pmatrix}-1&-2&-2&1\\1&-1&1&1\\1&-2&-2&1\\0&-2&2&-1\end{pmatrix}$, encuentre $M_{2,3}$ y $C_{2,3}$.

\item Para $A=\begin{pmatrix}-2&-2&1&2\\-1&-1&1&-2\\-2&-2&2&0\\0&0&0&-1\end{pmatrix}$, encuentre $M_{4,4}$ y $C_{4,4}$.

\end{enumerate}
\end{Ejer}

\begin{Ejer}

\textbf{Calcule la matriz de cofactores, transpuesta y, si el determinante es distinto de cero, la inversa. Compare con Gauss--Jordan.}
\begin{enumerate}[start=187]

\item Calcule $\adj(A)$ y $A^{-1}$ usando cofactores, con $A=\begin{pmatrix}1&0\\0&1\end{pmatrix}$.

\item Calcule $\adj(A)$ y $A^{-1}$ usando cofactores, con $A=\begin{pmatrix}1&0\\0&1\end{pmatrix}$.

\item Calcule $\adj(A)$ y $A^{-1}$ usando cofactores, con $A=\begin{pmatrix}3&-1\\-2&1\end{pmatrix}$.

\item Calcule $\adj(A)$ y $A^{-1}$ usando cofactores, con $A=\begin{pmatrix}1&1\\1&2\end{pmatrix}$.

\item Calcule $\adj(A)$ y $A^{-1}$ usando cofactores, con $A=\begin{pmatrix}3&2\\1&1\end{pmatrix}$.

\item Calcule $\adj(A)$ y $A^{-1}$ usando cofactores, con $A=\begin{pmatrix}1&1\\-1&0\end{pmatrix}$.

\item Calcule $\adj(A)$ y $A^{-1}$ usando cofactores, con $A=\begin{pmatrix}1&0&0\\0&-1&3\\0&-1&2\end{pmatrix}$.

\item Calcule $\adj(A)$ y $A^{-1}$ usando cofactores, con $A=\begin{pmatrix}1&-1&4\\0&1&0\\0&0&1\end{pmatrix}$.

\item Calcule $\adj(A)$ y $A^{-1}$ usando cofactores, con $A=\begin{pmatrix}1&0&0\\-1&1&-2\\2&0&1\end{pmatrix}$.

\item Calcule $\adj(A)$ y $A^{-1}$ usando cofactores, con $A=\begin{pmatrix}1&0&-1\\2&1&-1\\0&0&1\end{pmatrix}$.

\end{enumerate}
\end{Ejer}

\begin{Ejer}

\textbf{Reduzca cada matriz por filas y cuente pivotes. Los tamaños incluyen $2\times5$, $3\times4$, $4\times5$ y $5\times4$.}
\begin{enumerate}[start=197]

\item Determine $\rang(A)$ reduciendo $A=\begin{pmatrix}2&0&2\\-1&0&-1\end{pmatrix}$.

\item Determine $\rang(A)$ reduciendo $A=\begin{pmatrix}1&1&4&3\\-1&0&-2&-1\end{pmatrix}$.

\item Determine $\rang(A)$ reduciendo $A=\begin{pmatrix}1&2\\0&0\\0&0\end{pmatrix}$.

\item Determine $\rang(A)$ reduciendo $A=\begin{pmatrix}1&0&0&1\\-1&1&0&-2\\-1&1&0&-2\end{pmatrix}$.

\item Determine $\rang(A)$ reduciendo $A=\begin{pmatrix}1&0&0\\0&1&0\\0&1&1\\0&-1&-1\end{pmatrix}$.

\item Determine $\rang(A)$ reduciendo $A=\begin{pmatrix}1&0&-2&-1&0\\0&1&-1&1&-1\\0&0&0&0&0\\0&1&-1&1&-1\end{pmatrix}$.

\item Determine $\rang(A)$ reduciendo $A=\begin{pmatrix}1&0&0&1\\0&1&0&-1\\0&1&1&0\\0&1&1&0\\0&0&0&0\end{pmatrix}$.

\item Determine $\rang(A)$ reduciendo $A=\begin{pmatrix}1&0&0&0&-1\\1&1&1&1&-1\\0&0&1&1&-2\end{pmatrix}$.

\item Determine $\rang(A)$ reduciendo $A=\begin{pmatrix}1&0\\0&1\\0&0\\0&0\\0&0\end{pmatrix}$.

\item Determine $\rang(A)$ reduciendo $A=\begin{pmatrix}1&1&-1&1&-1\\0&0&0&0&0\end{pmatrix}$.

\item Determine $\rang(A)$ reduciendo $A=\begin{pmatrix}1&1&0&3\\0&1&0&1\\0&1&1&2\\0&0&0&0\end{pmatrix}$.

\item Determine $\rang(A)$ reduciendo $A=\begin{pmatrix}1&0&0&1&3\\0&1&0&0&-1\\0&0&1&1&3\\0&0&0&1&1\\0&0&0&0&0\end{pmatrix}$.

\item Determine $\rang(A)$ reduciendo $A=\begin{pmatrix}1&2&0\\0&0&0\\0&0&0\\0&0&0\\0&0&0\end{pmatrix}$.

\item Determine $\rang(A)$ reduciendo $A=\begin{pmatrix}1&0&-1\\0&1&1\\1&0&-1\end{pmatrix}$.

\item Determine $\rang(A)$ reduciendo $A=\begin{pmatrix}1&-1&0&0&-1\\0&1&0&0&2\\0&0&1&0&-1\\-1&0&0&1&-1\end{pmatrix}$.

\item Determine $\rang(A)$ reduciendo $A=\begin{pmatrix}1&0&0&2\\0&1&-2&2\\0&0&0&0\\0&0&0&0\\0&0&0&0\end{pmatrix}$.
\end{enumerate}
\end{Ejer}

\begin{Ejer}

\textbf{Resuelva $A\vect x=0$, indique las variables libres y escriba el núcleo en forma paramétrica. Verifique cada vector generador multiplicando por $A$.}
\begin{enumerate}[start=213]

\item Determine $\rang(A)$ y $\Ker(A)$ para $A=\begin{pmatrix}1&2&0\\-2&-4&0\end{pmatrix}$.

\item Determine $\rang(A)$ y $\Ker(A)$ para $A=\begin{pmatrix}1&1&0&2\\1&2&-1&2\end{pmatrix}$.

\item Determine $\rang(A)$ y $\Ker(A)$ para $A=\begin{pmatrix}1&0\\0&1\\-2&0\end{pmatrix}$.

\item Determine $\rang(A)$ y $\Ker(A)$ para $A=\begin{pmatrix}1&0&-2&0\\1&1&-4&2\\1&1&-4&2\end{pmatrix}$.

\item Determine $\rang(A)$ y $\Ker(A)$ para $A=\begin{pmatrix}1&0&0\\0&1&0\\-1&0&1\\0&-1&0\end{pmatrix}$.

\item Determine $\rang(A)$ y $\Ker(A)$ para $A=\begin{pmatrix}1&0&-2&-1&-2\\1&1&-4&0&-1\\0&0&0&0&0\\-1&0&2&1&2\end{pmatrix}$.

\item Determine $\rang(A)$ y $\Ker(A)$ para $A=\begin{pmatrix}1&1&0&1\\0&1&0&1\\0&0&1&-1\\0&0&0&0\\1&1&0&1\end{pmatrix}$.

\item Determine $\rang(A)$ y $\Ker(A)$ para $A=\begin{pmatrix}1&0&-1&3&0\\0&1&0&-1&0\\0&-1&1&0&2\end{pmatrix}$.

\item Determine $\rang(A)$ y $\Ker(A)$ para $A=\begin{pmatrix}1&0\\1&1\\1&1\\0&0\\0&0\end{pmatrix}$.

\item Determine $\rang(A)$ y $\Ker(A)$ para $A=\begin{pmatrix}1&2&2&2&2\\0&0&0&0&0\end{pmatrix}$.

\item Determine $\rang(A)$ y $\Ker(A)$ para $A=\begin{pmatrix}1&1&0&1\\0&1&0&-1\\0&0&1&2\\0&-1&0&1\end{pmatrix}$.

\item Determine $\rang(A)$ y $\Ker(A)$ para $A=\begin{pmatrix}1&-1&0&0&2\\0&1&0&0&-2\\0&0&1&0&1\\0&1&0&1&0\\0&0&0&0&0\end{pmatrix}$.

\item Determine $\rang(A)$ y $\Ker(A)$ para $A=\begin{pmatrix}1&-2&-1\\0&0&0\\1&-2&-1\\0&0&0\\0&0&0\end{pmatrix}$.

\item Determine $\rang(A)$ y $\Ker(A)$ para $A=\begin{pmatrix}1&0&0\\-1&1&-1\\-1&0&0\end{pmatrix}$.

\item Determine $\rang(A)$ y $\Ker(A)$ para $A=\begin{pmatrix}1&0&0&0&2\\0&1&0&-2&3\\0&0&1&0&2\\0&0&0&1&-2\end{pmatrix}$.

\item Determine $\rang(A)$ y $\Ker(A)$ para $A=\begin{pmatrix}1&0&-1&-2\\0&1&0&-1\\0&0&0&0\\0&0&0&0\\0&0&0&0\end{pmatrix}$.
\end{enumerate}
\end{Ejer}



\section*{Bibliografía de referencia}
\begin{itemize}
\item Grossman, S. I. (2012). \emph{Álgebra Lineal}. 7a edición. McGraw-Hill.
\item Lay, D. C. (2013). \emph{Álgebra lineal para cursos con enfoque por competencias}. Pearson.
\item Kolman, B. (2013). \emph{Álgebra Lineal}. Pearson Educación.
\item Poole, D. (2011). \emph{Álgebra lineal: una introducción moderna}. 3a edición. Cengage Learning.
\end{itemize}

\vspace{0.8em}
\end{document}
